\oldpage[IV-3]{95}
\section{Préschémas de Jacobson}
\label{section:IV.10-fr}

On a déjà eu l'occasion d'observer (5.2.5) que même les préschémas
\emph{excellents} (7.8.5) ne se comportent pas toujours comme les
«~variétés~» de la Géométrie algébrique classique, notamment en ce qui
concerne les questions de dimension~; c'est ainsi que si $X$ est le spectre
d'un anneau de valuation discrète complet, l'ensemble des points fermés de
$X$ (réduit à un seul élément) n'est pas partout dense dans $X$, et que son
complémentaire (réduit à un seul élément) est un ouvert partout dense dans
$X$ mais dont la dimension est nulle, donc $<\dim(X)$. On examine dans ce
paragraphe des conditions générales moyennant lesquelles de tels phénomènes
ne se présentent pas~; il en résulte un comportement plus satisfaisant à
certains égards pour les relations entre dimension, codimension, profondeur
et coprofondeur dans ces préschémas (10.6 et 10.8). En outre, dans les
préschémas considérés, le fait que l'ensemble des points fermés soit partout
dense (et même «~très dense~» (10.1.3)) permet de ne considérer que ces points
dans beaucoup de démonstrations~; on rejoint ainsi le point de vue classique
des «~variétés algébriques~» qui, de notre point de vue, sont les ensembles
de points fermés des \emph{préschémas algébriques sur un corps}, et on fait le
lien entre le langage des schémas et celui des «~variétés~» ou «~espaces
algébriques~» de Serre (10.9 et 10.10).

\subsection{Parties très denses d'un espace topologique}
\label{subsection:IV.10.1-fr}

\begin{env}[10.1.1]
\label{IV.10.1.1-fr}
On dit qu'une partie d'un espace topologique $X$ est
\emph{quasi-constructible} si elle est réunion finie de parties localement
fermées de $X$. On dit qu'une partie $T$ de $X$ est \emph{localement
quasi-constructible} si pour tout $x\in X$, il existe un voisinage ouvert $V$
de $x$ tel que $T\cap V$ soit quasi-constructible dans $V$. Il est clair que
toute partie quasi-constructible de $X$ est localement quasi-constructible~:
la réciproque est vraie si $X$ est quasi-compact. Le raisonnement de
($\mathbf 0_{\mathrm{III}}$, 9.1.3), où l'on supprime le mot
«~rétrocompact~», montre que l'ensemble des parties quasi-constructibles
(resp. localement quasi-constructibles) de $X$, que nous désignerons par
$\mathfrak{Qc}(X)$ (resp. $\mathfrak{Lqc}(X)$) est stable par réunion finie,
intersection finie et passage aux complémentaires. Si $f:X\to Y$ est une
application continue, il résulte aussitôt de ces définitions que, pour toute
partie quasi-constructible (resp. localement quasi-constructible) $Z$ de $Y$,
$f^{-1}(Z)$ est quasi-constructible (resp. localement quasi-constructible)
dans $X$.

Les parties constructibles (resp. localement constructibles) de $X$ sont
évidemment quasi-constructibles (resp. localement quasi-constructibles)~; la
réciproque est vraie lorsque $X$ est noethérien (resp. localement
noethérien).

Dans ce qui suit, nous désignerons respectivement par $\mathfrak O(X)$,
$\mathfrak F(X)$, $\mathfrak{Lf}(X)$, $\mathfrak C(X)$, $\mathfrak{Lc}(X)$
l'ensemble des parties de $X$ qui sont respectivement ouvertes, fermées,
localement fermées, constructibles, localement constructibles.
\end{env}

\oldpage[IV-3]{96}
\begin{proposition}[10.1.2]
\label{IV.10.1.2-fr}
Soient $X$ un espace topologique, $X_0$ une partie de $X$. Les conditions
suivantes sont équivalentes~:
\begin{enumerate}[label=\emph{\alph*)}]
  \item Pour toute partie localement fermée $Z\neq\varnothing$ de $X$, on a
  $Z\cap X_0\neq\varnothing$.
  \item[\emph{a$'$)}] Pour toute partie fermée $Z$ de $X$, on a
  $Z=\overline{Z\cap X_0}$.
  \item Pour toute partie $Z\neq\varnothing$ de $X$ localement
  quasi-constructible, on a $Z\cap X_0\neq\varnothing$.
  \item[\emph{b$'$)}] Pour toute partie localement quasi-constructible $Z$ de
  $X$, on a $Z\subset\overline{Z\cap X_0}$ (autrement dit $Z\cap X_0$ est
  \emph{dense} dans $Z$).
  \item L'application $U\mapsto U\cap X_0$ de $\mathfrak O(X)$ dans
  $\mathfrak O(X_0)$ est injective (donc bijective).
  \item[\emph{c$'$)}] L'application $F\mapsto F\cap X_0$ de $\mathfrak F(X)$
  dans $\mathfrak F(X_0)$ est injective (donc bijective).
  \item[\emph{c$''$)}] L'application $Z\mapsto Z\cap X_0$ de
  $\mathfrak{Qc}(X)$ dans $\mathfrak{Qc}(X_0)$ est injective (donc bijective).
  \item[\emph{c$'''$)}] L'application $Z\mapsto Z\cap X_0$ de
  $\mathfrak{Lqc}(X)$ dans $\mathfrak{Lqc}(X_0)$ est injective.
\end{enumerate}
En outre, lorsque ces conditions sont vérifiées, l'application
$Z\mapsto Z\cap X_0$ de $\mathfrak{Lqc}(X)$ dans $\mathfrak{Lqc}(X_0)$ est
bijective.
\end{proposition}

Notons que les assertions de surjectivité dans \emph{c)} et \emph{c$'$)} sont
triviales~: elles entraînent que toute partie localement fermée dans $X_0$ est
trace sur $X_0$ d'une partie localement fermée dans $X$, donc l'application
$\mathfrak{Qc}(X)\to\mathfrak{Qc}(X_0)$ définie dans \emph{c$''$)} est aussi
surjective.

Nous prouverons les implications
\[
c''')\Rightarrow c'')\Rightarrow c')\Rightarrow c)\Rightarrow b)
\Rightarrow b')\Rightarrow a')\Rightarrow a)\Rightarrow c''').
\]
Les trois premières sont triviales. Pour voir que \emph{c)} implique
\emph{b)}, notons d'abord que \emph{c)} entraîne que $X_0$ est \emph{dense}
dans $X$. Remplaçant $X$ et $X_0$ par un ouvert convenable $U$ de $X$ et par
$U\cap X_0$ respectivement, on peut donc se borner au cas où $Z$ est
localement fermée dans $X$, donc $Z=V\cap\complement W$, où $V$ et $W$ sont
ouverts dans $X$~; l'hypothèse $Z\neq\varnothing$ signifie que
$V\not\subset W$, ou encore $V\cup W\neq W$. En vertu de \emph{c)}, on a donc
$(V\cup W)\cap X_0\neq W\cap X_0$, donc $V\cap X_0\not\subset W$, et par
suite $(V\cap\complement W)\cap X_0\neq\varnothing$.

Pour voir que \emph{b)} entraîne \emph{b$'$)}, il suffit d'appliquer
\emph{b)} à $Z\cap U$, où $U$ est un voisinage ouvert arbitraire d'un point de
$Z$. Comme $Z\cap\overline{X_0}\subset\overline{Z\cap X_0}$, il est trivial
que \emph{b$'$)} implique \emph{a$'$)}. Pour montrer que \emph{a$'$)} entraîne
\emph{a)}, remarquons que si $Z$ est localement fermée dans $X$, on peut
écrire $Z=F-F'$ où $F$ et $F'$ sont fermés dans $X$ et $F'\subset F$~; d'où
$Z\cap X_0=(F\cap X_0)-(F'\cap X_0)$. Si l'on avait
$Z\cap X_0=\varnothing$, on en déduirait $F\cap X_0=F'\cap X_0$, donc
$F=F'$ en vertu de \emph{a$'$)}, c'est-à-dire $Z=\varnothing$.

Pour voir que \emph{a)} entraîne \emph{c$'''$)}, il suffit de montrer que si
$Z\neq\varnothing$ est localement quasi-constructible, on a
$Z\cap X_0\neq\varnothing$~: en effet, la relation
$Z'\cap X_0=Z''\cap X_0$ est équivalente à
$((Z'\cup Z'')-(Z'\cap Z''))\cap X_0=\varnothing$. Autrement dit, il suffit de
prouver que \emph{a)} entraîne \emph{b)}~; en outre, remplaçant $X$ et $X_0$
par un ouvert $U$ de $X$ et par $U\cap X_0$ respectivement, on est bien
ramené au cas où $Z$ est localement fermée dans $X$, d'où la conclusion.

Reste à montrer que l'application
$\mathfrak{Lqc}(X)\to\mathfrak{Lqc}(X_0)$ est surjective. Soit donc $Z_0$ une
partie localement quasi-constructible de $X_0$~: il y a un recouvrement
$(V_\alpha)$ de $X_0$ par des ouverts de $X_0$ tel que $Z_0\cap V_\alpha$ soit
quasi-constructible dans $V_\alpha$ (et
\oldpage[IV-3]{97}
donc aussi dans $X_0$). En vertu de \emph{c)}, il y a pour tout $\alpha$ un
seul ouvert $U_\alpha$ de $X$ tel que $X_0\cap U_\alpha=V_\alpha$, et en
vertu de \emph{c$''$)} un seul ensemble quasi-constructible $Z_\alpha$ dans
$X$ tel que $Z_0\cap V_\alpha=X_0\cap Z_\alpha$. Si $\alpha$ et $\beta$ sont
deux indices quelconques, on a
\[
Z_\beta\cap U_\alpha\cap X_0
=U_\alpha\cap Z_0\cap V_\beta
=V_\alpha\cap Z_0\cap V_\beta
=Z_\alpha\cap X_0\cap V_\beta
\subset Z_\alpha\cap X_0~;
\]
comme $Z_\beta\cap U_\alpha$ et $Z_\alpha$ sont quasi-constructibles dans
$X$, il résulte de \emph{c$''$)} que
$Z_\beta\cap U_\alpha\subset Z_\alpha$. Si l'on pose
$Z=\bigcup_\alpha Z_\alpha$, on a donc $Z\cap U_\alpha=Z_\alpha$ pour tout
$\alpha$~; d'ailleurs, comme les $V_\alpha$ recouvrent $X_0$, il résulte de
\emph{c)} que les $U_\alpha$ recouvrent $X$, et l'on voit donc que $Z$ est
localement quasi-constructible dans $X$ et $Z_0=Z\cap X_0$.

\begin{definition}[10.1.3]
\label{IV.10.1.3-fr}
Lorsqu'une partie $X_0$ d'un espace topologique vérifie les conditions
équivalentes de (10.1.2), on dit que $X_0$ est \emph{très dense} dans $X$.
\end{definition}

On a déjà vu au cours de la démonstration de (10.1.2) que $X_0$ est alors
\emph{dense} dans $X$.

\begin{corollary}[10.1.4]
\label{IV.10.1.4-fr}
Si $X_0$ est très dense dans $X$ et $U$ une partie ouverte de $X$,
$U\cap X_0$ est très dense dans $U$. Inversement, si $(U_\alpha)$ est un
recouvrement ouvert de $X$ tel que $U_\alpha\cap X_0$ soit très dense dans
$U_\alpha$ pour tout $\alpha$, alors $X_0$ est très dense dans $X$.
\end{corollary}

Comme toute partie localement fermée dans $U$ est localement fermée dans $X$,
la première assertion résulte du critère \emph{a)} de (10.1.2)~; il en est de
même de la seconde, car si $Z\neq\varnothing$ est localement fermée dans $X$,
$Z\cap U_\alpha$ est localement fermée dans $U_\alpha$ pour tout $\alpha$,
et $Z\cap U_\alpha\neq\varnothing$ pour un $\alpha$ au moins.

\subsection{Quasi-homéomorphismes}
\label{subsection:IV.10.2-fr}

\begin{proposition}[10.2.1]
\label{IV.10.2.1-fr}
Soient $X_0$, $X$ deux espaces topologiques, $f:X_0\to X$ une application
continue. Les conditions suivantes sont équivalentes~:
\begin{enumerate}[label=\emph{\alph*)}]
  \item L'application $U\mapsto f^{-1}(U)$ de $\mathfrak O(X)$ dans
  $\mathfrak O(X_0)$ est bijective.
  \item[\emph{a$'$)}] L'application $F\mapsto f^{-1}(F)$ de $\mathfrak F(X)$
  dans $\mathfrak F(X_0)$ est bijective.
  \item La topologie de $X_0$ est image réciproque par $f$ de celle de $X$, et
  $f(X_0)$ est très dense (10.1.3) dans $X$.
  \item Le foncteur $\mathcal F\mapsto f^*(\mathcal F)$ de la catégorie des
  faisceaux d'ensembles (resp. des faisceaux de groupes abéliens) sur $X$ dans
  la catégorie des faisceaux d'ensembles (resp. des faisceaux de groupes
  abéliens) sur $X_0$ est une équivalence de catégories.
\end{enumerate}
\end{proposition}

Il est clair que \emph{a)} et \emph{a$'$)} sont équivalentes, et \emph{a)}
implique que la topologie de $X_0$ est image réciproque par $f$ de celle de
$X$. D'autre part, si $f(X_0)$ n'est pas très dense dans $X$, il y a deux
ouverts distincts $U_1$, $U_2$ de $X$ tels que
$U_1\cap f(X_0)=U_2\cap f(X_0)$, et par suite
$f^{-1}(U_1)=f^{-1}(U_2)$, ce qui achève de montrer que \emph{a)} entraîne
\emph{b)}. Inversement, la condition \emph{b)} implique que les applications
$U\mapsto U\cap f(X_0)$ de $\mathfrak O(X)$ dans $\mathfrak O(f(X_0))$ et
$V\mapsto f^{-1}(V)$ de $\mathfrak O(f(X_0))$ dans $\mathfrak O(X_0)$ sont
bijectives, donc il en est de même de leur composée $U\mapsto f^{-1}(U)$.

Pour voir que \emph{a)} entraîne \emph{c)}, il suffit d'appliquer la définition
de $f^*(\mathcal F)$ ($\mathbf 0_{\mathrm I}$, 3.7.1) et les axiomes des
faisceaux~: \emph{a)} entraîne que pour tout ouvert $U$ de $X$, l'application
canonique
$\Gamma(U,\mathcal F)\to\Gamma(f^{-1}(U),f^*(\mathcal F))$ est une bijection
fonctorielle en $\mathcal F$, d'où \emph{c)}. Il reste à montrer que
\emph{c)} entraîne \emph{b)}.

\oldpage[IV-3]{98}
Supposons d'abord que $f(X_0)$ ne soit pas très dense dans $X$~; il existe
alors deux parties fermées distinctes $Y\supset Y'$ de $X$ telles que
$Y\cap f(X_0)=Y'\cap f(X_0)$. Soit $\mathcal F$ (resp. $\mathcal F'$) le
faisceau de groupes abéliens sur $X$ image directe par l'injection canonique
$Y\to X$ (resp. $Y'\to X$) du faisceau simple associé au préfaisceau constant
$\mathbf Z$ sur $Y$ (resp. $Y'$)~; la définition de $f^*$
($\mathbf 0_{\mathrm I}$, 3.7.1) montre que $f^*(\mathcal F)$ est isomorphe à
$f^*(\mathcal F')$ mais $\mathcal F$ n'est pas isomorphe à $\mathcal F'$,
donc la condition \emph{c)} n'est pas vérifiée. (On notera que le foncteur
$f^*$ n'est même pas alors \emph{fidèle}, car si $u$ et $v$ sont
l'automorphisme identique et l'endomorphisme nul de $\mathcal F/\mathcal F'$,
$f^*(u)$ et $f^*(v)$ sont égaux.)

Montrons maintenant que \emph{c)} entraîne que la topologie de $X_0$ est
l'image réciproque de celle de $X$ par $f$. Notons que si la condition
\emph{c)} est vérifiée pour la catégorie des faisceaux d'ensembles, elle l'est
aussi pour la catégorie des faisceaux de groupes abéliens, car il résulte
aussitôt des définitions ($\mathbf 0_{\mathrm{III}}$, 8.2.5) que cette dernière
n'est autre que la catégorie des \emph{objets en groupes abéliens} dans la
catégorie des faisceaux d'ensembles. Il suffira donc de prouver notre assertion
en supposant \emph{c)} vérifiée pour les catégories de faisceaux de groupes
abéliens sur $X$ et $X_0$. Or, si $\mathbf Z_X$ désigne le faisceau simple sur
$X$ associé au préfaisceau constant égal à $\mathbf Z$, on a un isomorphisme
canonique
\[
\Gamma(X,\mathcal F)\xrightarrow{\sim}\operatorname{Hom}(\mathbf Z_X,\mathcal F)
\]
fonctoriel en $\mathcal F$ (même raisonnement que dans
($\mathbf 0_{\mathrm I}$, 5.1.1))~; comme il est clair par définition
($\mathbf 0_{\mathrm I}$, 3.7.1) que $f^*(\mathbf Z_X)=\mathbf Z_{X_0}$,
l'hypothèse que $f^*$ est une équivalence entraîne que l'homomorphisme
canonique $\Gamma(X,\mathcal F)\to\Gamma(X_0,f^*(\mathcal F))$ est
\emph{bijectif} pour tout faisceau de groupes abéliens $\mathcal F$ sur $X$.
Or, soit $Y_0$ une partie fermée de $X_0$~; soit $\mathcal F_0$ le faisceau
sur $X$ image directe par l'injection canonique $Y_0\to X_0$ du faisceau
simple associé au préfaisceau constant $\mathbf Z$ sur $Y_0$. Comme $f^*$ est
une équivalence, $\mathcal F_0$ est isomorphe à un faisceau de la forme
$f^*(\mathcal F)$, où $\mathcal F$ est un faisceau de groupes abéliens sur
$X$. Considérons alors la section $s_0$ de $\mathcal F_0$ au-dessus de $X_0$
telle que $(s_0)_{x_0}=1_{x_0}$ si $x_0\in Y_0$,
$(s_0)_{x_0}=0_{x_0}$ si $x_0\notin Y_0$. Les remarques précédentes montrent
qu'il existe une section $s$ de $\mathcal F$ au-dessus de $X$ telle que
$(s_0)_{x_0}=s_{f(x_0)}$ pour tout $x_0\in X_0$ (les fibres
$(\mathcal F_0)_{x_0}$ et $\mathcal F_{f(x_0)}$ étant canoniquement
identifiées ($\mathbf 0_{\mathrm I}$, 3.7.1))~; cela entraîne que $Y_0$ est
l'image réciproque par $f$ de l'ensemble $Y$ des $x\in X$ tels que
$s_x\neq0_x$, et $y$ est fermé dans $X$. C.Q.F.D.

\begin{definition}[10.2.2]
\label{IV.10.2.2-fr}
Lorsqu'une application continue $f:X_0\to X$ vérifie les conditions
équivalentes de (10.2.1), on dit que $f$ est un \emph{quasi-homéomorphisme} de
$X_0$ dans $X$.
\end{definition}

En vertu de (10.2.1, \emph{b)}), dire qu'une partie $X_0$ d'un espace
topologique $X$ est \emph{très dense} dans $X$ signifie donc que l'injection
canonique $X_0\to X$ est un \emph{quasi-homéomorphisme}.

\begin{corollary}[10.2.3]
\label{IV.10.2.3-fr}
Le composé de deux quasi-homéomorphismes est un quasi-homéomorphisme.
\end{corollary}

Cela résulte aussitôt de (10.2.1, \emph{a)}).

\begin{corollary}[10.2.4]
\label{IV.10.2.4-fr}
Soient $f:X\to Y$ un quasi-homéomorphisme, $Y'$ une partie localement
quasi-constructible de $Y$, $X'=f^{-1}(Y')$~; alors la restriction
$f'=f|X':X'\to Y'$ est un quasi-homéomorphisme.
\end{corollary}

Il est clair en vertu de (10.2.1, \emph{b)}) que la topologie induite sur $X'$
par celle de $X$ est l'image réciproque par $f'$ de la topologie induite sur
$Y'$ par celle de $Y$. D'autre part,
\oldpage[IV-3]{99}
soit $Z\neq\varnothing$ une partie fermée \emph{dans} $Y'$, et $z\in Z$~; il
y a un voisinage ouvert $U$ de $z$ dans $Y$ tel que $U\cap Y'$ soit réunion
d'un nombre fini de parties $Y'_j$ fermées dans $U$~; si $j$ est un indice tel
que $z\in Y'_j$, $Z\cap Y'_j$ est donc fermé dans $U$. Puisque $f(X)\cap U$
est très dense dans $U$ (10.1.4), $Z\cap Y'_j\cap f(X)$ n'est pas vide
(10.1.2, \emph{a)}), ni \emph{a fortiori} $Z\cap f(X)$~; mais comme
$Z\subset Y'$, on a $Z\cap f(X)=Z\cap f'(X')$~; le critère
(10.1.2, \emph{a)}) montre donc que $f'(X')$ est très dense dans $Y'$, et
l'on conclut à l'aide de (10.2.1, \emph{b)}).

\begin{corollary}[10.2.5]
\label{IV.10.2.5-fr}
Soient $f:X\to Y$ une application continue, $(V_\alpha)$ un recouvrement
ouvert de $Y$. Si, pour tout $\alpha$, la restriction
$f^{-1}(V_\alpha)\to V_\alpha$ de $f$ est un quasi-homéomorphisme, alors $f$
est un quasi-homéomorphisme.
\end{corollary}

Cela résulte aussitôt du critère (10.2.1, \emph{b)}) et de (10.1.4).

\begin{corollary}[10.2.6]
\label{IV.10.2.6-fr}
Soient $f:X\to Y$ un quasi-homéomorphisme, $Y'$ une partie localement
quasi-constructible de $Y$, $X'=f^{-1}(Y')$. Pour que $Y'$ soit quasi-compact
(resp. noethérien, resp. rétrocompact dans $Y$), il faut et il suffit que $X'$
soit quasi-compact (resp. noethérien, resp. rétrocompact dans $X$).
\end{corollary}

Démontrons d'abord les deux premières assertions~; en vertu de (10.2.4), on
peut se borner au cas où $Y'=Y$. Dire que $X$ est quasi-compact (resp.
noethérien) signifie que pour toute famille filtrante croissante $(U_\alpha)$
dans $\mathfrak O(X)$ ayant $X$ pour plus grand élément (resp. pour toute
famille filtrante croissante $(U_\alpha)$ dans $\mathfrak O(X)$), il existe
$\gamma$ tel que $U_\alpha=U_\gamma$ pour $\alpha\geq\gamma$. Comme
$U\mapsto f^{-1}(U)$ est une bijection de $\mathfrak O(Y)$ sur
$\mathfrak O(X)$, notre assertion résulte aussitôt de la remarque précédente.

Les ouverts quasi-compacts de $X$ sont donc les ensembles de la forme
$f^{-1}(U)$ où $U$ est ouvert quasi-compact dans $Y$, en vertu de
(10.2.1, \emph{a)}) et de ce qui précède. Pour que $X'$ soit rétrocompact dans
$X$, il faut et il suffit alors que pour tout ouvert quasi-compact $U$ dans
$Y$, $f^{-1}(U)\cap X'=f^{-1}(U\cap Y')$ soit quasi-compact
($\mathbf 0_{\mathrm{III}}$, 9.1.1)~; la première partie de la démonstration
montre que cela équivaut à dire que $U\cap Y'$ est quasi-compact pour tout
ouvert quasi-compact $U$, c'est-à-dire que $Y'$ est rétrocompact dans $Y$.

\begin{proposition}[10.2.7]
\label{IV.10.2.7-fr}
Soit $f:X\to Y$ un quasi-homéomorphisme. Alors l'application
$Z\mapsto f^{-1}(Z)$ de $\mathfrak P(Y)$ dans $\mathfrak P(X)$ définit par
restriction les bijections suivantes (cf. (10.1.1) pour les notations)~:
\[
\begin{gathered}
\mathfrak O(Y)\xrightarrow{\sim}\mathfrak O(X)\\
\mathfrak F(Y)\xrightarrow{\sim}\mathfrak F(X)\\
\mathfrak{Lf}(Y)\xrightarrow{\sim}\mathfrak{Lf}(X)\\
\mathfrak{Qc}(Y)\xrightarrow{\sim}\mathfrak{Qc}(X)\\
\mathfrak{Lqc}(Y)\xrightarrow{\sim}\mathfrak{Lqc}(X)\\
\mathfrak C(Y)\xrightarrow{\sim}\mathfrak C(X)\\
\mathfrak{Lc}(Y)\xrightarrow{\sim}\mathfrak{Lc}(X).
\end{gathered}
\]
\end{proposition}

\oldpage[IV-3]{100}
Pour les deux premières, ce n'est autre que la définition (10.2.2)~; comme la
topologie de $X$ est l'image réciproque par $f$ de celle de $f(X)$ les cinq
dernières applications, où l'on remplace $Y$ par $f(X)$, sont bijectives. On
peut donc (par (10.2.1, \emph{b)})) se borner au cas où $X$ est un sous-espace
très dense de $Y$, et le fait que les applications
$\mathfrak{Lf}(Y)\to\mathfrak{Lf}(X)$,
$\mathfrak{Qc}(Y)\to\mathfrak{Qc}(X)$,
$\mathfrak{Lqc}(Y)\to\mathfrak{Lqc}(X)$ sont bijectives a déjà été prouvé
(10.1.2). Toute partie localement constructible étant localement
quasi-constructible, les applications $\mathfrak C(Y)\to\mathfrak C(X)$ et
$\mathfrak{Lc}(Y)\to\mathfrak{Lc}(X)$ sont injectives~; en outre, pour tout
ouvert $U\subset Y$, la restriction $f^{-1}(U)\to U$ de $f$ est un
quasi-homéomorphisme (10.2.4), donc si l'on montre que
$\mathfrak C(Y)\to\mathfrak C(X)$ est surjective, il en sera de même de
$\mathfrak{Lc}(Y)\to\mathfrak{Lc}(X)$. Mais en vertu de (10.2.6), toute partie
ouverte rétrocompacte $Z$ dans $X$ est de la forme $f^{-1}(T)$, où $T$ est
ouvert rétrocompact dans $Y$~; cela prouve évidemment la surjectivité de
$\mathfrak C(Y)\to\mathfrak C(X)$.

\begin{remarks}[10.2.8]
\label{IV.10.2.8-fr}
\upshape (i) On a vu dans la démonstration de (10.2.1) que si $f$ est un
quasi-homéomorphisme, l'application canonique
\begin{equation}
\label{IV.10.2.8.1-fr}
\tag{10.2.8.1}
\Gamma(Y,\mathcal F)\to\Gamma(X,f^*(\mathcal F))
\end{equation}
est un isomorphisme de groupes abéliens fonctoriel en $\mathcal F$ dans la
catégorie des faisceaux de groupes abéliens sur $Y$. Comme $f^*$ est exact
dans cette catégorie ($\mathbf 0_{\mathrm I}$, 3.7.2), cela implique que les
homomorphismes canoniques (T, 3.2.2)
\[
\mathrm H^i(Y,\mathcal F)\xrightarrow{\sim}
\mathrm H^i(X,f^*(\mathcal F))
\]
sont \emph{bijectifs} pour tout $i$.

(ii) Si $f:X\to Y$ est une application continue et $\mathcal F$,
$\mathcal G$ deux faisceaux de groupes abéliens sur $Y$, on a
\[
f^*(\mathcal F\otimes_{\mathbf Z_Y}\mathcal G)
=f^*(\mathcal F)\otimes_{\mathbf Z_X}f^*(\mathcal G)
\]
à un isomorphisme canonique près ($\mathbf 0_{\mathrm I}$, 4.3.3). On en
conclut que si $f$ est un quasi-homéomorphisme, $f^*$ est encore une
\emph{équivalence} de la catégorie des faisceaux d'anneaux sur $Y$, et de la
catégorie des faisceaux d'anneaux sur $X$~; la donnée d'une structure d'espace
annelé sur $Y$ est donc équivalente à celle d'une structure d'espace annelé
sur $X$.

Étant donnés deux espaces annelés $(X,\mathcal O_X)$, $(Y,\mathcal O_Y)$, nous
dirons qu'un morphisme d'espaces annelés
$f=(\psi,\theta):(X,\mathcal O_X)\to(Y,\mathcal O_Y)$ est un
\emph{quasi-isomorphisme} si $\psi$ est un quasi-homéomorphisme de $X$ dans
$Y$ et $\theta^\sharp:\psi^*(\mathcal O_Y)\to\mathcal O_X$ un isomorphisme de
faisceaux d'anneaux~; lorsqu'il en est ainsi, l'espace annelé
$(X,\mathcal O_X)$ est entièrement déterminé, à isomorphisme près, par
$(Y,\mathcal O_Y)$, l'espace $X$ et le quasi-homéomorphisme $\psi$ (que l'on
peut prendre arbitraire). Lorsque $f$ est un quasi-isomorphisme d'espaces
annelés, le foncteur
\[
\mathcal F\mapsto f^*(\mathcal F)
\]
est une \emph{équivalence} de la catégorie des $\mathcal O_Y$-Modules et de
celle des $\mathcal O_X$-Modules, puisque $f^*(\mathcal F)$ s'identifie
canoniquement ici à $\psi^*(\mathcal F)$. On en conclut par exemple des
isomorphismes de bi-$\partial$-foncteurs
\[
\operatorname{Ext}^i_{\mathcal O_Y}(\mathcal F,\mathcal G)
\xrightarrow{\sim}
\operatorname{Ext}^i_{\mathcal O_X}(f^*(\mathcal F),f^*(\mathcal G)).
\]
De façon générale, on peut dire que les constructions usuelles de la théorie
des faisceaux et de l'algèbre homologique, faites sur l'espace annelé $Y$ ou
l'espace annelé $X$, sont équivalentes.
\end{remarks}

\oldpage[IV-3]{101}
\subsection{Espaces de Jacobson}
\label{subsection:IV.10.3-fr}

\begin{definition}[10.3.1]
\label{IV.10.3.1-fr}
On dit qu'un espace topologique $X$ est un \emph{espace de Jacobson} si
l'ensemble $X_0$ des points fermés de $X$ est très dense dans $X$ (autrement
dit, si l'injection canonique $X_0\to X$ est un
quasi-homéomorphisme).
\end{definition}

Cela signifie donc (10.1.2) que toute partie localement fermée (ou seulement
localement quasi-constructible) $Z\neq\varnothing$ de $X$ contient un point
fermé de $X$, ou encore que \emph{toute partie fermée de $X$ est l'adhérence
de l'ensemble de ses points fermés}.

\begin{proposition}[10.3.2]
\label{IV.10.3.2-fr}
Soient $X$ un espace de Jacobson, $Z$ une partie localement
quasi-constructible de $X$~; alors le sous-espace $Z$ de $X$ est un espace de
Jacobson, et pour qu'un point de $Z$ soit fermé dans $Z$, il faut et il suffit
qu'il soit fermé dans $X$.
\end{proposition}

Si $X_0$ est l'ensemble des points fermés de $X$, $Z\cap X_0$ est très dense
dans $Z$ en vertu de (10.2.4) appliqué à l'injection $i:X_0\to X$~; comme
l'ensemble $Z_0$ des points fermés de $Z$ contient évidemment $Z\cap X_0$,
il est \emph{a fortiori} très dense dans $Z$, donc $Z$ est un espace de
Jacobson. Montrons maintenant qu'en fait on a $Z_0=Z\cap X_0$~; soit $x\in Z$
un point fermé dans $Z$~; soit $\overline{\{x\}}$ son adhérence dans $X$~;
$\overline{\{x\}}\cap Z=\{x\}$ est donc une partie localement
quasi-constructible de $X$, et comme son intersection avec $X_0$ est non vide
(10.1.2), on a $x\in X_0$.

\begin{proposition}[10.3.3]
\label{IV.10.3.3-fr}
Soient $X$ un espace topologique, $(U_\alpha)$ un recouvrement ouvert de $X$.
Pour que $X$ soit un espace de Jacobson, il faut et il suffit que chacun des
sous-espaces $U_\alpha$ le soit.
\end{proposition}

La condition est nécessaire en vertu de (10.3.2). Inversement, montrons
d'abord que l'hypothèse que les $U_\alpha$ sont des espaces de Jacobson
entraîne que pour qu'un point $x\in U_\alpha$ soit fermé dans $X$, il suffit
qu'il soit fermé dans $U_\alpha$. Il suffit en effet de voir que cette
condition entraîne que $x$ est aussi fermé dans chacun des $U_\beta$ qui le
contiennent~; mais $U_\alpha\cap U_\beta$ est ouvert dans $U_\alpha$, donc $x$
est fermé dans $U_\alpha\cap U_\beta$, et en vertu de (10.3.2), $x$ est aussi
fermé dans $U_\beta$, ce qui achève la démonstration.

\subsection{Préschémas de Jacobson et anneaux de Jacobson}
\label{subsection:IV.10.4-fr}

\begin{definition}[10.4.1]
\label{IV.10.4.1-fr}
On dit qu'un préschéma $X$ est un \emph{préschéma de Jacobson} si l'espace
topologique $X$ sous-jacent est un espace de Jacobson. On dit qu'un anneau
$A$ est un \emph{anneau de Jacobson} si $\operatorname{Spec}(A)$ est un
préschéma de Jacobson.
\end{definition}

Toute partie fermée de $X=\operatorname{Spec}(A)$ est de la forme
$Z=V(\mathfrak a)$, où $\mathfrak a$ est un idéal égal à sa racine, et
l'ensemble $X_0$ des points fermés de $X$ est l'ensemble des idéaux maximaux
de $A$~; dire que $Z\cap X_0$ est dense dans $Z$ signifie donc que
$\mathfrak a$ est \emph{intersection d'idéaux maximaux} (\textbf{I}, 1.1.4)~;
comme $\mathfrak a$ est intersection d'idéaux premiers, il revient au même de
dire que tout idéal premier de $A$ est intersection d'idéaux maximaux~; en
vertu de (10.3.1) et (10.1.2), la définition usuelle des anneaux de Jacobson
(Bourbaki, \emph{Alg. comm.}, chap.~V, \S~3, n\textsuperscript{o}~4, déf.~1)
coïncide donc avec la définition (10.4.1).

\oldpage[IV-3]{102}
\begin{proposition}[10.4.2]
\label{IV.10.4.2-fr}
Soient $X$ un préschéma, $(U_\alpha)$ un recouvrement de $X$ formé d'ouverts
affines. Pour que $X$ soit un préschéma de Jacobson, il faut et il suffit que
les anneaux $\Gamma(U_\alpha,\mathcal O_X)$ des $U_\alpha$ soient des anneaux
de Jacobson.
\end{proposition}

Cela résulte de (10.3.3) et de la définition (10.4.1).

\begin{env}[10.4.3]
\label{IV.10.4.3-fr}
Un préschéma discret est un préschéma de Jacobson~; un anneau artinien est donc
un anneau de Jacobson. Tout anneau principal ayant une infinité d'idéaux
maximaux (par exemple $\mathbf Z$) est un anneau de Jacobson~; un anneau local
noethérien $A$ est un anneau de Jacobson si et seulement si son idéal maximal
est son seul idéal premier, c'est-à-dire si $A$ est artinien. Tout
sous-préschéma d'un préschéma de Jacobson est un préschéma de Jacobson en vertu
de (10.3.2).
\end{env}

\begin{proposition}[10.4.4]
\label{IV.10.4.4-fr}
Soit $B$ un anneau intègre. Les conditions suivantes sont équivalentes~:
\begin{enumerate}[label=\emph{\alph*)}]
  \item Il existe $f\neq0$ dans $B$ tel que $B_f$ soit un corps.
  \item Le corps des fractions de $B$ est une $B$-algèbre de type fini.
  \item Il existe un corps $K$ contenant $B$ et qui est une $B$-algèbre de
  type fini.
  \item Le point générique de $\operatorname{Spec}(B)$ est isolé dans
  $\operatorname{Spec}(B)$.
\end{enumerate}
\end{proposition}

Il est clair que \emph{d)} est équivalent à \emph{a)}, puisque \emph{d)}
signifie qu'il existe $f\neq0$ dans $B$ tel que $D(f)$ soit réduit au point
générique de $\operatorname{Spec}(B)$. Il est trivial que \emph{a)} entraîne
\emph{b)} et que \emph{b)} entraîne \emph{c)}. Enfin, \emph{c)} entraîne
\emph{a)}, en vertu de Bourbaki, \emph{Alg. comm.}, chap.~V, \S~3,
n\textsuperscript{o}~1, cor.~2 du th.~1.

\begin{proposition}[10.4.5]
\label{IV.10.4.5-fr}
Soit $A$ un anneau. Les conditions suivantes sont équivalentes~:
\begin{enumerate}[label=\emph{\alph*)}]
  \item $A$ est un anneau de Jacobson.
  \item Pour tout idéal premier non maximal $\mathfrak p$ de $A$ et tout
  $f\neq0$ dans $B=A/\mathfrak p$, $B_f$ n'est pas un corps.
  \item[\emph{b$'$)}] Toute $A$-algèbre de type fini $K$ qui est un corps est
  une $A$-algèbre finie.
\end{enumerate}
\end{proposition}

On sait que \emph{a)} entraîne \emph{b$'$)} (Bourbaki,
\emph{Alg. comm.}, chap.~V, \S~3, n\textsuperscript{o}~4, cor.~3 du th.~3).
En outre, le noyau de l'homomorphisme $A\to K$ est alors un idéal
\emph{maximal} $\mathfrak m$ de $A$, et $K$ est une extension \emph{finie} de
$A/\mathfrak m$ (\emph{loc. cit.}). Il est trivial que \emph{b$'$)} entraîne
\emph{b)}, puisque $A/\mathfrak p=B$ n'est pas un corps, toute $B$-algèbre de
type fini est une $A$-algèbre de type fini et $B_f$ est une $B$-algèbre de
type fini. Reste à voir que \emph{b)} implique \emph{a)}. Nous utiliserons le
lemme suivant~:

\begin{lemma}[10.4.5.1]
\label{IV.10.4.5.1-fr}
Soit $X$ un espace topologique ayant la propriété suivante~: pour toute partie
localement fermée $Z\neq\varnothing$ de $X$, il existe une partie $Z'$ de $Z$,
localement fermée dans $X$ (ou dans $Z$, ce qui revient au même), et un point
$x\in Z'$, fermé dans $Z'$. Alors, pour que $X$ soit un espace de Jacobson, il
faut et il suffit qu'aucun point non fermé $x$ de $X$ ne soit isolé dans
$\overline{\{x\}}$.
\end{lemma}

Si $X$ est un espace de Jacobson, un point non fermé $x$ de $X$ ne peut être
isolé dans $\overline{\{x\}}$, car cela signifierait qu'il existe un ouvert
$U$ de $X$ tel que $U\cap\overline{\{x\}}=\{x\}$~; or
$U\cap\overline{\{x\}}$ est localement fermé dans $X$ et ne contiendrait aucun
point fermé dans $X$, ce qui est contraire à l'hypothèse que $X$ est un espace
de Jacobson. Inversement, supposons vérifiée la condition de l'énoncé~; alors
l'ensemble des points fermés de $X$ est identique à l'ensemble $X_0$ des
$x\in X$ qui sont isolés dans $\overline{\{x\}}$~; mais il résulte de
(5.1.10.1) que

\oldpage[IV-3]{103}
cet ensemble est \emph{très dense} dans $X$, donc $X$ est un espace de
Jacobson par définition.

Rappelons (5.1.10) que l'hypothèse faite sur $X$ dans (10.4.5.1) est toujours
vérifiée lorsque $X$ est l'espace sous-jacent à un \emph{préschéma}.

Revenant alors à la démonstration de (10.4.5), la condition \emph{b)} entraîne
que pour tout point $x$ non fermé de $\operatorname{Spec}(A)$, le point
générique $x$ de
$\operatorname{Spec}(A/\mathfrak j_x)=V(\mathfrak j_x)=\overline{\{x\}}$
n'est pas isolé dans $\overline{\{x\}}$, donc \emph{b)} entraîne \emph{a)} en
vertu du lemme (10.4.5.1) et de (10.4.4).

\begin{corollary}[10.4.6]
\label{IV.10.4.6-fr}
Toute algèbre de type fini $B$ sur un anneau de Jacobson $A$ est un anneau de
Jacobson, et l'image réciproque dans $A$ de tout idéal maximal de $B$ est un
idéal maximal de $A$. En particulier, toute algèbre de type fini sur un corps
ou sur $\mathbf Z$ est un anneau de Jacobson.
\end{corollary}

Une $B$-algèbre de type fini $K$ qui est un corps est aussi une $A$-algèbre de
type fini, donc un $A$-module de type fini et \emph{a fortiori} un $B$-module
de type fini, d'où la première assertion~; la seconde a été démontrée au cours
de la démonstration de (10.4.5), appliqué à $K=B/\mathfrak m$.

\begin{corollary}[10.4.7]
\label{IV.10.4.7-fr}
Si $X$ est un préschéma de Jacobson, $f:Y\to X$ un morphisme localement de type
fini, alors $Y$ est un préschéma de Jacobson et l'image par $f$ de tout point
fermé dans $Y$ est un point fermé dans $X$.
\end{corollary}

La question étant locale sur $X$ et sur $Y$, on est ramené au cas où $X$ et
$Y$ sont affines, et le corollaire résulte alors de (10.4.6).

\begin{corollary}[10.4.8]
\label{IV.10.4.8-fr}
Si $k$ est un corps algébriquement clos, $X$ un $k$-préschéma localement de
type fini, l'ensemble des points de $X$ rationnels sur $k$ est très dense dans
$X$.
\end{corollary}

En effet, $X$ est un préschéma de Jacobson (10.4.7) et les points fermés de
$X$ sont exactement les points rationnels sur $k$ (\textbf{I}, 6.4.2).

\begin{env}[10.4.9]
\label{IV.10.4.9-fr}
Le fait que les préschémas localement de type fini sur un corps ou sur
$\mathbf Z$ sont des préschémas de Jacobson est particulièrement important,
en raison de la possibilité de se ramener à ce cas dans de nombreuses
questions de Géométrie algébrique (8.1.2, \emph{c)}). Nous allons en donner
deux exemples~:
\end{env}

\begin{env}[10.4.10]
\label{IV.10.4.10-fr}
\textbf{Applications (10.4.10)~: I~: Démonstration de (6.15.9).} --- Soient
$k$ un corps séparablement clos, $X$ un $k$-préschéma localement de type fini
sur $k$ et \emph{unibranche}. On sait que la fermeture intégrale d'une
$k$-algèbre intègre de type fini $A$ dans une extension finie de son corps des
fractions est une $A$-algèbre finie (Bourbaki, \emph{Alg. comm.}, chap.~V,
\S~3, n\textsuperscript{o}~2, th.~2), donc toute $k$-algèbre de type fini est
un anneau universellement japonais~; on en conclut que l'ensemble des points
$x\in X$ où $X$ est géométriquement unibranche est localement constructible
(9.7.10). Mais l'hypothèse et le lemme (6.15.8) entraînent que cet ensemble
contient tous les points \emph{fermés} de $X$. La conclusion résulte donc de
(10.4.6), (10.3.1) et de la bijectivité de l'application canonique
$\mathfrak{Lc}(X)\to\mathfrak{Lc}(X_0)$ (où $X_0$ est l'ensemble des points
fermés de $X$) (10.2.7).
\end{env}

\begin{env}[10.4.11]
\label{IV.10.4.11-fr}
\textbf{Applications~: II~: Proposition (10.4.11).} --- Soient $S$ un
préschéma, $X$ un $S$-préschéma de type fini. Tout $S$-endomorphisme de $X$
qui est radiciel est surjectif (donc bijectif).
\end{env}

\oldpage[IV-3]{104}
Soient $f:X\to S$ le morphisme structural, $g:X\to X$ le $S$-endomorphisme
considéré et, pour tout $s\in S$, soit $g_s$ le morphisme déduit de $g$ par le
changement de base $\operatorname{Spec}(k(s))\to S$, qui est un
$k(s)$-endomorphisme de la fibre
$f^{-1}(s)=X\times_S\operatorname{Spec}(k(s))$.

Pour prouver que $g$ est surjectif, il suffit de prouver que $g_s$ est
surjectif pour tout $s\in S$, donc on peut (en vertu de (\textbf{I}, 3.5.7))
se borner au cas où $S=\operatorname{Spec}(k)$ est le spectre d'un corps $k$,
auquel cas $f$ est un morphisme de présentation finie, puisque $S$ est
noethérien. Appliquant (8.9.1) et (8.10.5, (vii)), on est ramené au cas où
$S=\operatorname{Spec}(A)$, où $A$ est une sous-$\mathbf Z$-algèbre de type
fini de $k$. Or $X$ est alors un préschéma de Jacobson (10.4.7), et $g(X)$ est
constructible dans $X$ (\textbf{I}, 1.8.5), donc, pour prouver que $g(X)=X$,
il suffit de montrer que $g(X)$ contient tous les points \emph{fermés} de $X$
(10.3.1).

\begin{lemma}[10.4.11.1]
\label{IV.10.4.11.1-fr}
Soit $Y$ un $\mathbf Z$-préschéma de type fini.
\begin{enumerate}[label=(\roman*)]
  \item Pour qu'un point $y\in Y$ soit fermé dans $Y$, il faut et il suffit
  que $k(y)$ soit un corps fini.
  \item Pour tout nombre premier $p$ et tout entier $d\geq1$, l'ensemble des
  points $y\in Y$ tels que $k(y)$ soit une extension de $\mathbf F_p$ dont le
  degré divise $d$ est fini.
\end{enumerate}
\end{lemma}

L'assertion (i) résulte de ce que l'image d'un point fermé $y\in Y$ dans
$\operatorname{Spec}(\mathbf Z)$ est un point fermé (10.4.7), autrement dit
un nombre premier, et de (\textbf{I}, 6.4.2). D'autre part, comme $Y$ est
réunion finie d'ouverts affines de type fini sur $\mathbf Z$, on peut se
borner pour prouver (ii) au cas où $Y=\operatorname{Spec}(C)$, où $C$ est une
$\mathbf Z$-algèbre de type fini. Or, les points $y\in Y$ tels que le degré de
$k(y)$ sur $\mathbf F_p$ divise $d$ correspondent biunivoquement aux
homomorphismes $C\to\mathbf F_{p^d}$~; mais si $(t_i)$ est un système de
générateurs de la $\mathbf Z$-algèbre $C$, tout homomorphisme de $C$ est
déterminé par ses valeurs aux éléments $t_i$, et par suite il n'y a qu'un
nombre fini d'homomorphismes de $C$ dans un corps fini.

Cela étant, $X$ est un $\mathbf Z$-préschéma de type fini~; soit $T_{p,d}$
l'ensemble des points fermés $z\in X$ tels que $k(z)$ soit une extension de
$\mathbf F_p$ de degré divisant $d$~; il résulte de (10.4.11.1) que l'ensemble
$T_{p,d}$ est fini et que l'ensemble des points fermés de $X$ est réunion des
$T_{p,d}$. En outre, si $z\in T_{p,d}$ et si $h$ est un endomorphisme
quelconque de $X$, $k(h(z))$ est isomorphe à un sous-corps de $k(z)$, donc
$h(z)\in T_{p,d}$, autrement dit $T_{p,d}$ est \emph{stable} par tout
endomorphisme de $X$. Comme $g$ est par hypothèse injectif, sa restriction à
$T_{p,d}$ est \emph{bijective} puisque $T_{p,d}$ est fini, ce qui achève de
démontrer la proposition.

Nous verrons plus loin (17.9.7) que lorsqu'on suppose de plus, d'une part que
$X$ est un $S$-préschéma de présentation finie, d'autre part que $g$ est un
monomorphisme, alors on peut affirmer que $g$ est un \emph{automorphisme} de
$X$.

\subsection{Préschémas de Jacobson noethériens}
\label{subsection:IV.10.5-fr}

\begin{proposition}[10.5.1]
\label{IV.10.5.1-fr}
Soit $B$ un anneau intègre noethérien. Les conditions équivalentes \emph{a)} à
\emph{d)} de (10.4.4) sont alors aussi équivalentes aux suivantes~:
\begin{enumerate}[label=\emph{\alph*)}]
  \setcounter{enumi}{4}
  \item $\operatorname{Spec}(B)$ est fini.
  \item $B$ est un anneau semi-local de dimension $\leq1$.
\end{enumerate}
\end{proposition}

Il résulte en effet du théorème d'Artin--Tate ($\mathbf 0$, 16.3.3) que les
conditions \emph{a)} et \emph{f)} sont équivalentes. La condition \emph{f)}
implique que $\operatorname{Spec}(B)$ est réunion de l'ensemble

\oldpage[IV-3]{105}
fini de ses points fermés et de son point générique, donc \emph{f)} implique
\emph{e)} sans supposer $A$ noethérien. Enfin \emph{e)} implique \emph{d)} sans
supposer $A$ noethérien, car le point générique $x$ de $X$ est le complémentaire
de la réunion des adhérences $\overline{\{y\}}$, où $y$ parcourt l'ensemble des
points $y\neq x$, et comme ces points sont en nombre fini, $\{x\}$ est
complémentaire d'un ensemble fermé dans $X$.

\begin{corollary}[10.5.2]
\label{IV.10.5.2-fr}
Soit $A$ un anneau noethérien. Pour que $A$ soit un anneau de Jacobson, il faut
et il suffit qu'il n'existe aucun idéal premier $\mathfrak p$ de $A$ tel que
$A/\mathfrak p$ soit un anneau semi-local de dimension $1$.
\end{corollary}

Cela résulte aussitôt de (10.5.1) et de la condition \emph{b)} de (10.4.5),
les idéaux premiers $\mathfrak p$ de $A$ tels que $A/\mathfrak p$ soit
semi-local de dimension $0$ étant les idéaux maximaux de $A$.

\begin{corollary}[10.5.3]
\label{IV.10.5.3-fr}
Soit $X$ un préschéma noethérien irréductible. Les conditions suivantes sont
équivalentes~:
\begin{enumerate}[label=\emph{\alph*)}]
  \item Le point générique de $X$ est isolé.
  \item $X$ est fini.
  \item $X$ est de dimension $\leq1$ et l'ensemble de ses points fermés est
  fini.
\end{enumerate}
\end{corollary}

Il existe par hypothèse un nombre fini d'ouverts affines irréductibles $U_i$
($1\leq i\leq n$) recouvrant $X$, et dont chacun contient donc le point
générique de $X$~; il suffit de prouver l'équivalence de \emph{a)}, \emph{b)}
et \emph{c)} pour chacun des $(U_i)_{\mathrm{red}}$ (compte tenu de
($\mathbf 0$, 14.1.7)). Mais cette équivalence résulte alors de (10.5.1).

\begin{remark}[10.5.4]
\label{IV.10.5.4-fr}
Un préschéma noethérien $X$ vérifiant les conditions équivalentes de (10.5.3)
n'est pas nécessairement un schéma affine~; en fait il peut même être
\emph{non séparé}. On en a un exemple en remplaçant, dans l'exemple
(\textbf{I}, 5.5.11) de «~droite affine avec point dédoublé~», $X_1$ et $X_2$
par le spectre de l'anneau de valuation discrète $(K[s])_{(s)}$, $U_{12}$ et
$U_{21}$ par l'ouvert réduit au point générique dans $X_1$ et $X_2$
respectivement~; le préschéma non séparé $X$ que l'on obtient a exactement
$3$ points.
\end{remark}

\begin{proposition}[10.5.5]
\label{IV.10.5.5-fr}
Soit $X$ un préschéma noethérien.
\begin{enumerate}[label=(\roman*)]
  \item L'ensemble $X_0$ des points $x\in X$ tels que $\overline{\{x\}}$ soit
  fini est très dense dans $X$.
  \item Pour que $X$ soit un préschéma de Jacobson, il faut et il suffit qu'il
  n'existe aucun sous-préschéma de $X$ isomorphe au spectre d'un anneau
  semi-local intègre de dimension $1$.
\end{enumerate}
\end{proposition}

La condition que $\overline{\{x\}}$ soit fini équivaut en effet ici (10.5.3)
au fait que $x$ soit isolé dans $\overline{\{x\}}$, $\overline{\{x\}}$ étant
espace sous-jacent d'un sous-préschéma (noethérien) de $X$~; l'assertion (i)
résulte donc de (5.1.10.1). De même, compte tenu de (10.4.5.1), pour démontrer
l'assertion (ii), remarquons que pour qu'un point non fermé $x$ de $X$
appartienne à $X_0$, il faut et il suffit que le sous-préschéma fermé intègre
$Y$ de $X$ ayant $\overline{\{x\}}$ pour espace sous-jacent soit de dimension
$1$ et fini, donc réunion finie de sous-préschémas ouverts (dans $Y$) affines
$U_i$ qui sont des spectres d'anneaux semi-locaux intègres de dimension $1$.
Inversement, s'il y a un sous-préschéma $Z$ de $X$ qui soit spectre d'un anneau
semi-local intègre de dimension $1$, $Z$ n'est pas un préschéma de Jacobson
(10.5.2), donc il en est de même de $X$ (10.3.2).

\oldpage[IV-3]{106}
\begin{remark}[10.5.6]
\label{IV.10.5.6-fr}
L'assertion (ii) de (10.5.5) est encore valable lorsque $X$ est
\emph{localement noethérien}~: en effet, si $(V_\alpha)$ est un recouvrement de
$X$ formé d'ouverts affines (noethériens), tout sous-préschéma d'un $V_\alpha$
est un sous-préschéma de $X$~; inversement, si un sous-préschéma $Z$ de $X$ est
isomorphe au spectre d'un anneau semi-local intègre de dimension $1$, il y a
un $\alpha$ tel que $V_\alpha\cap Z$ contienne un ouvert affine $U$ de $Z$ non
réduit au point générique de $Z$, et qui est donc aussi spectre d'un anneau
semi-local intègre de dimension $1$. On conclut à l'aide de (10.3.3).
\end{remark}

\begin{proposition}[10.5.7]
\label{IV.10.5.7-fr}
Soient $X$ un préschéma localement noethérien, $Y$ une partie fermée de $X$
telle que toute partie fermée non vide de $X$ rencontre $Y$. Alors le
préschéma induit sur l'ouvert $X-Y$ est un préschéma de Jacobson.
\end{proposition}

Appliquons le critère (10.5.5, (ii)), et supposons qu'il y ait un
sous-préschéma $Z$ de $X-Y$ qui soit spectre d'un anneau semi-local intègre de
dimension $1$, le point générique $z$ de $Z$ étant donc isolé dans $Z$ (ou
dans l'adhérence $\overline Z$ de $Z$ dans $X$), et $Z$ étant distinct de
$\{z\}$. Soit $y\neq z$ un point de $Z$~; comme il n'appartient pas à $Y$, il
n'est pas fermé dans $X$, et son adhérence $\overline{\{y\}}$ dans $X$
rencontre $Y$ en un point $x\neq y$ qui est donc spécialisation de $y$.
L'existence de la chaîne
$\{x\}\subset\overline{\{y\}}\subset\overline{\{z\}}$ montre alors que la
dimension de $\overline{\{z\}}$ serait $\geq2$, et il en serait de même de la
dimension de $U\cap\overline{\{z\}}$, où $U$ est un voisinage ouvert affine
(donc noethérien) de $x$ dans $X$. Or, cela contredit le fait que le point
générique $z$ de $U\cap\overline{\{z\}}$ est isolé dans
$U\cap\overline{\{z\}}$ (10.5.3).

\begin{corollary}[10.5.8]
\label{IV.10.5.8-fr}
Soit $A$ un anneau noethérien~; pour tout élément $f$ du radical $\mathfrak R$
de $A$, l'anneau $A_f$ est un anneau de Jacobson, et l'ouvert
$\operatorname{Spec}(A)-V(\mathfrak R)$ est un schéma de Jacobson.
\end{corollary}

Si $X=\operatorname{Spec}(A)$, $Y=V(\mathfrak J)$, où $\mathfrak J$ est un
idéal de $A$, dire que toute partie fermée non vide de $X$ rencontre $Y$
équivaut ($\mathbf 0_{\mathrm I}$, 2.1.3) à dire que $Y$ contient tous les
\emph{points fermés} de $X$, ou encore que $\mathfrak J$ est \emph{contenu
dans le radical} $\mathfrak R$ de $A$. Si $f\in\mathfrak R$, l'ouvert $D(f)$
ne rencontre donc pas $V(\mathfrak R)$, et est un espace de Jacobson en vertu
de (10.5.7).

\begin{corollary}[10.5.9]
\label{IV.10.5.9-fr}
Soient $A$ un anneau local noethérien, $\mathfrak m$ son idéal maximal,
$X=\operatorname{Spec}(A)$, $Y=X-\{\mathfrak m\}$~; alors $Y$ est un schéma de
Jacobson, dont les points fermés sont les idéaux premiers
$\mathfrak p\in\operatorname{Spec}(A)$ tels que $\dim(A/\mathfrak p)=1$.
\end{corollary}

La première assertion est un cas particulier de (10.5.8)~; d'autre part les
points fermés de $Y$ sont les idéaux premiers $\mathfrak p$ de $A$ qui sont
des éléments maximaux dans l'ensemble des idéaux premiers $\neq\mathfrak m$,
ce qui, par définition de la dimension, signifie que
$\dim(A/\mathfrak p)=1$.

\begin{proposition}[10.5.10]
\label{IV.10.5.10-fr}
Soit $A$ un anneau local noethérien réduit et complet et qui n'est pas un
corps. Alors les intersections finies des noyaux des homomorphismes locaux de
$A$ dans des anneaux de valuation discrète $V$, qui font de $V$ une
$A$-algèbre finie, forment une base de filtre tendant vers $0$ pour la
topologie adique de $A$.
\end{proposition}

Il suffit (Bourbaki, \emph{Alg. comm.}, chap.~III, \S~2,
n\textsuperscript{o}~7, prop.~8) de prouver que l'intersection des noyaux
envisagés dans l'énoncé est réduite à $0$. Supposons d'abord que $A$ soit
intègre et de dimension $1$~; en vertu du théorème de Nagata ($\mathbf 0$,
23.1.5 et 23.1.6), la clôture intégrale $A'$ de $A$ est un anneau local
complet intègre et

\oldpage[IV-3]{107}
intégralement clos et de dimension $1$ et une $A$-algèbre finie~; c'est donc
un anneau de valuation discrète, et la proposition en résulte aussitôt dans
ce cas.

Passons au cas général~; soit $x$ le point fermé de
$X=\operatorname{Spec}(A)$, et posons $Y=X-\{x\}$~; on sait (10.5.9) que $Y$
est un préschéma de Jacobson. Montrons que cela entraîne que l'intersection
des idéaux premiers $\mathfrak p$ de $A$ tels que
$\dim(A/\mathfrak p)=1$ est \emph{réduite à $0$}~; la proposition en résultera
puisque, pour chacun de ces idéaux $\mathfrak p$, l'intersection des noyaux des
homomorphismes locaux de $A/\mathfrak p$ dans des anneaux de valuation
discrète $V$, faisant de $V$ une $(A/\mathfrak p)$-algèbre finie, est réduite
à $0$. Mais dire que l'intersection de ces idéaux premiers est réduite à $0$
signifie que l'ensemble de ces idéaux est dense dans $X$, ou encore dans $Y$
(puisque $Y$ est dense dans $X$), et cela résulte aussitôt de (10.5.9).

\subsection{Dimension dans les préschémas de Jacobson}
\label{subsection:IV.10.6-fr}

Les résultats de ce numéro précisent dans certains cas et généralisent des résultats
du \S~5.

\begin{proposition}[10.6.1]
\label{IV.10.6.1-fr}
Soit $S$ un préschéma localement noethérien, vérifiant en outre les conditions
suivantes~: $1^\circ$ $S$ est un préschéma de Jacobson~; $2^\circ$ pour tout
$s\in S$, $\mathcal O_{S,s}$ est universellement caténaire (5.6.2)~; $3^\circ$ toute
composante irréductible $S'$ de $S$ est équicodimensionnelle (autrement dit, pour
tout point fermé $s$ de $S'$ et tout sous-préschéma de $S$ ayant $S'$ pour espace
sous-jacent, on a $\dim(S')=\dim(\mathcal O_{S',s})$). On a alors les propriétés
suivantes~:
\begin{enumerate}[label=(\roman*)]
  \item Pour tout morphisme localement de type fini $g:X\to S$, $X$ vérifie les
  conditions $1^\circ$, $2^\circ$ et $3^\circ$ précédentes. En particulier, si $X$
  est équidimensionnel (par exemple si $X$ est irréductible), $X$ est
  biéquidimensionnel (autrement dit, $X$ est caténaire et pour tout point fermé
  $x$ de $X$, on a $\dim(\mathcal O_{X,x})=\dim(X)$ ($\mathbf 0$, 14.3.3)).
  \item Soient $X$, $Y$ deux $S$-préschémas localement de type fini sur $S$,
  $f:X\to Y$ un $S$-morphisme~; on suppose $X$ irréductible et $f$ dominant. Si
  $\xi$ (resp. $\eta$) est le point générique de $X$ (resp. $Y$) et
  $e=\dim(f^{-1}(\eta))=\deg.\operatorname{tr}_{k(\eta)}k(\xi)$, on a
  \begin{equation}
  \label{IV.10.6.1.1-fr}
  \tag{10.6.1.1}
  \dim(X)=\dim(Y)+e.
  \end{equation}
  \item Soient $X$, $Y$ deux $S$-préschémas localement de type fini sur $S$,
  $f:X\to Y$ un $S$-morphisme, $n$ un entier $\geq0$ tel que l'on ait
  $\dim(f^{-1}(y))\geq n$ (resp. $\dim(f^{-1}(y))\leq n$) pour tout $y\in Y$.
  Alors on a
  \begin{equation}
  \label{IV.10.6.1.2-fr}
  \tag{10.6.1.2}
  \dim(X)\geq\dim(Y)+n
  \end{equation}
  (resp.
  \begin{equation}
  \label{IV.10.6.1.3-fr}
  \tag{10.6.1.3}
  \dim(X)\leq\dim(Y)+n).
  \end{equation}
\end{enumerate}
\end{proposition}

(i) La propriété $1^\circ$ pour $X$ résulte de (10.4.7). Pour tout $x\in X$,
$\mathcal O_{X,x}$ est anneau local d'une $\mathcal O_{S,g(x)}$-algèbre de type fini en
un idéal premier, et l'homomorphisme
$\mathcal O_{S,g(x)}\to\mathcal O_{X,x}$

\oldpage[IV-3]{108}
est local~; donc (5.6.3, (iv)) $\mathcal O_{X,x}$ est universellement caténaire. Pour
démontrer que $X$ vérifie la condition $3^\circ$, considérons plusieurs cas~:

\emph{a)} $X$ est un sous-préschéma fermé irréductible de $S$~; soient $S'$ une
composante irréductible de $S$ contenant $X$, $\xi$ le point générique de $X$,
$x$ un point fermé de $X$~; pour tout $s\in S'$, $\mathcal O_{S',s}$, quotient de
$\mathcal O_{S,s}$, est caténaire (5.6.1), donc les conditions $2^\circ$ et $3^\circ$
entraînent que $S'$ est biéquidimensionnel ($\mathbf 0$, 14.3.3). En vertu de
(5.1.2) et de ($\mathbf 0$, 14.3.3.2), on a donc
\[
  \dim(\mathcal O_{X,x})=\dim(\mathcal O_{S',x})-\dim(\mathcal O_{S',\xi})
  =\dim(S')-\dim(\mathcal O_{S',\xi})
\]
ce qui prouve que $\dim(\mathcal O_{X,x})$ ne dépend pas du point fermé $x$
considéré, d'où l'assertion dans ce cas (5.1.4).

\emph{b)} $X$ est irréductible et $g$ dominant. Alors, pour tout point fermé $x$
de $X$, $g(x)$ est fermé dans $S$ par (10.4.7)~; comme $\mathcal O_{S,g(x)}$ est
universellement caténaire, il résulte de (5.6.5.3) que l'on a
\[
  \dim(\mathcal O_{X,x})=\dim(\mathcal O_{S,g(x)})+e
\]
où $e=\dim(g^{-1}(\zeta))$, $\zeta$ étant le point générique de $S$. Comme
$\dim(S)=\dim(\mathcal O_{S,g(x)})$ en vertu de la condition $3^\circ$ pour $S$, on
a $\dim(\mathcal O_{X,x})=\dim(S)+e$ pour tout point fermé $x\in X$~; cela prouve
la condition $3^\circ$ pour $X$ (5.1.4), et en même temps la formule
\begin{equation}
\label{IV.10.6.1.4-fr}
\tag{10.6.1.4}
\dim(X)=\dim(S)+e.
\end{equation}

\emph{c)} Cas général. En considérant un sous-préschéma réduit $X'$ de $X$ ayant
pour espace sous-jacent une composante irréductible de $X$, on est ramené au cas
où $X$ est intègre~; utilisant (\textbf{I}, 5.2.2) et le cas \emph{a)} prouvé
ci-dessus, on peut ensuite remplacer $S$ par le sous-préschéma réduit ayant
$\overline{g(X)}$ pour espace sous-jacent~; on est alors ramené au cas \emph{b)},
et cela achève de démontrer (i).

(ii) Le morphisme $f$ étant localement de type fini (\textbf{I}, 1.3.4), on peut
appliquer les résultats de (i) en remplaçant $S$ par $Y$~; en outre, comme $X$
est irréductible et $f$ dominant, on peut aussi remplacer $S$ par $Y$ dans
(10.6.1.4), ce qui donne (10.6.1.1).

(iii) L'assertion relative au cas où $\dim(f^{-1}(y))\leq n$ pour tout $y\in Y$
a déjà été démontrée sous des hypothèses plus générales dans (5.6.7). Supposons
que $\dim(f^{-1}(y))\geq n$ pour tout $y\in Y$, et considérons un point générique
$\eta$ d'une composante irréductible $Y'$ de $Y$~; il existe au moins une
composante irréductible $Z$ de $f^{-1}(\eta)$ de dimension $\geq n$~; si $\xi$
est le point générique de $Z$, $\xi$ est aussi point générique d'une composante
irréductible $X'$ de $X$ telle que $f(X')$ soit dense dans $Y'$ ($\mathbf 0_{\mathrm I}$,
2.1.8). Considérons les sous-préschémas réduits de $X$, $Y$ ayant respectivement
pour espaces sous-jacents $X'$, $Y'$, et la restriction $X'\to Y'$ de $f$
(\textbf{I}, 5.2.2)~; il résulte alors de (ii) que
$\dim(X')\geq\dim(Y')+n$, et a fortiori $\dim(X)\geq\dim(Y')+n$~; ceci étant vrai
pour toute composante irréductible $Y'$ de $Y$, on en conclut l'inégalité
(10.6.1.2).

\begin{corollary}[10.6.2]
\label{IV.10.6.2-fr}
Supposons que $X$ vérifie les conditions $1^\circ$, $2^\circ$ et $3^\circ$ de
(10.6.1). Alors, pour tout ouvert $U$ dense dans $X$, on a
$\dim(U)=\dim(X)$.
\end{corollary}

\oldpage[IV-3]{109}
On sait en effet que $\dim(U)\leq\dim(X)$ ($\mathbf 0$, 14.1.4)~; en outre,
comme $X$ est un préschéma de Jacobson, $U$ contient un point fermé de toute
composante irréductible de $X$, donc $\dim(U)=\dim(X)$ en vertu de (10.6.1) et
($\mathbf 0$, 14.1.2.1).

\begin{proposition}[10.6.3]
\label{IV.10.6.3-fr}
Supposons que $X$ vérifie les conditions $1^\circ$, $2^\circ$ et $3^\circ$ de
(10.6.1), et soit $Y$ une partie fermée de $X$. Alors, pour tout $x\in Y$, et
tout voisinage ouvert $U$ de $x$ dans $X$ ne rencontrant pas les composantes
irréductibles de $Y$ qui ne contiennent pas $x$, on a
\begin{equation}
\label{IV.10.6.3.1-fr}
\tag{10.6.3.1}
\dim_x(Y)=\dim(U\cap Y)=\dim(\overline{\{x\}})
 +\operatorname{codim}(\overline{\{x\}},Y)=\sup_i\dim(Y_i)
\end{equation}
où $Y_i$ ($1\leq i\leq m$) sont les composantes irréductibles de $Y$ contenant
$x$.
\end{proposition}

En considérant un sous-préschéma fermé de $X$ ayant $Y$ pour espace sous-jacent,
on peut se borner, en vertu de (10.6.1, (i)), au cas où $Y=X$. En vertu du choix
de $U$, $U$ est réunion des $U\cap X_i$, donc
$\dim(U)=\sup_i\dim(U\cap X_i)$~; mais d'après (10.6.2) appliqué à un
sous-préschéma fermé de $X$ ayant $X_i$ pour espace sous-jacent, on a (compte
tenu de (10.6.1, (i))) $\dim(U\cap X_i)=\dim(X_i)$~; cela démontre que le second
et le quatrième terme de (10.6.3.1) sont égaux. Comme $U\cap X_i$ est
biéquidimensionnel (10.6.1, (i)) (puisque l'immersion $U\cap X_i\to X_i$ est de
type fini (\textbf{I}, 6.3.5)), on a
\begin{equation}
\label{IV.10.6.3.2-fr}
\tag{10.6.3.2}
\dim(X_i)=\dim(U\cap X_i)=\dim(U\cap\overline{\{x\}})
 +\operatorname{codim}(U\cap\overline{\{x\}},U\cap X_i)
\end{equation}
($\mathbf 0$, 14.3.5). D'après (10.6.2) appliqué à un sous-préschéma fermé de
$X$ ayant $\overline{\{x\}}$ pour espace sous-jacent,
$\dim(U\cap\overline{\{x\}})=\dim(\overline{\{x\}})$~; comme on a aussi, pour les
mêmes raisons,
\begin{equation}
\label{IV.10.6.3.3-fr}
\tag{10.6.3.3}
\dim(X_i)=\dim(\overline{\{x\}})
 +\operatorname{codim}(\overline{\{x\}},X_i)
\end{equation}
on obtient
\[
  \dim(U)=\dim(\overline{\{x\}})
  +\sup_i\operatorname{codim}(\overline{\{x\}},X_i)
  =\dim(\overline{\{x\}})+\operatorname{codim}(\overline{\{x\}},X)
\]
par définition de la codimension ($\mathbf 0$, 14.2.1). Ceci montre que $\dim(U)$
est indépendante du voisinage ouvert $U$ de $x$ vérifiant les conditions de
l'énoncé, donc est égal à $\dim_x(X)$, par ($\mathbf 0$, 14.1.4.1).

\begin{corollary}[10.6.4]
\label{IV.10.6.4-fr}
Sous les hypothèses de (10.6.3), soient $\mathcal F$ un $\mathcal O_X$-Module cohérent,
$Y$ son support. Pour tout $x\in Y$, on a
\begin{equation}
\label{IV.10.6.4.1-fr}
\tag{10.6.4.1}
\dim(\overline{\{x\}})+\dim(\mathcal F_x)=\dim_xY.
\end{equation}
\end{corollary}

En effet, cela résulte de (10.6.3.1) et de la formule
$\dim(\mathcal F_x)=\operatorname{codim}(\overline{\{x\}},Y)$ (5.1.12.2).

\subsection{Exemples et contre-exemples}
\label{subsection:IV.10.7-fr}

\begin{env}[10.7.1]
\label{IV.10.7.1-fr}
Soit $S$ un préschéma localement noethérien de dimension $\leq1$ et supposons
que $S$ soit un \emph{préschéma de Jacobson}~; lorsque $S$ est noethérien, il
revient au même de dire que les composantes irréductibles de $S$ de dimension
$1$ sont infinies, car tout $x\in S$ qui n'est pas fermé est le point générique
d'une telle composante (10.4.5 et 10.5.4). Alors $S$ vérifie aussi les conditions
$2^\circ$ et $3^\circ$ de (10.6.1)~: en effet, tout anneau local $\mathcal O_{S,s}$
est de dimension $0$ ou $1$, et par suite est universellement caténaire (7.2.9)~;
d'autre part, une composante irréductible $S'$ de $S$ est, soit réduite

\oldpage[IV-3]{110}
à un point, soit de dimension $1$, et pour tout point fermé $s\in S'$,
$\mathcal O_{S',s}$ est nécessairement de dimension $1$.

On déduit de ces remarques et de (10.6.1) que tout préschéma localement de type
fini sur $S$ vérifie aussi les propriétés $1^\circ$, $2^\circ$ et $3^\circ$ de
(10.6.1)~: il en est ainsi en particulier des préschémas \emph{localement de type
fini sur un corps ou sur $\mathbf Z$}.
\end{env}

\begin{env}[10.7.2]
\label{IV.10.7.2-fr}
Soit $A$ un anneau local noethérien \emph{universellement caténaire} et soit $S$
le \emph{complémentaire dans $X=\operatorname{Spec}(A)$ du point fermé} $a$.
Alors $S$ vérifie les conditions $1^\circ$, $2^\circ$ et $3^\circ$ de (10.6.1)~:
en effet, on a déjà vu que $S$ est un préschéma de Jacobson (10.5.9)~; comme $A$
est universellement caténaire, il en est de même des anneaux locaux $A_{\mathfrak p}$
aux idéaux premiers de $A$ (5.6.3). D'autre part, une composante irréductible $S'$
de $S$ est le complémentaire de $a$ dans une composante irréductible $X'$ de $X$~;
pour tout point fermé $x$ de $S$, l'adhérence de $x$ dans $X$ est donc
$\{x,a\}$, autrement dit $\dim(\overline{\{x\}})=1$ et par suite, puisque $X'$
est par hypothèse biéquidimensionnel ($\mathbf 0$, 14.3.3), on a
$\dim(\mathcal O_{X',x})=\dim(X')-1$ ($\mathbf 0$, 14.3.2), ce qui prouve que $S'$
est équicodimensionnel (5.1.4).
\end{env}

\begin{env}[10.7.3]
\label{IV.10.7.3-fr}
Soit $A$ un anneau de valuation discrète~; montrons que dans l'anneau
$B=A[T_1,\ldots,T_n]$ des polynômes à $n$ indéterminées, il existe deux idéaux
maximaux $\mathfrak m$, $\mathfrak n$ de \emph{hauteurs respectives} $n$ et
$n+1$. On l'a vu dans (5.2.5, (i)) pour $n=1$~; démontrons-le par récurrence sur
$n$. Comme $A[T_1,\ldots,T_{n+1}]$ est un $A[T_1,\ldots,T_n]$-module libre,
donc fidèlement plat, il y a dans $A[T_1,\ldots,T_{n+1}]$ deux idéaux maximaux
$\mathfrak m'$, $\mathfrak n'$ respectivement au-dessus de $\mathfrak m$ et
$\mathfrak n$ ($\mathbf 0_{\mathrm I}$, 6.5.1)~; en outre, d'après (5.5.3), ces
idéaux sont nécessairement de hauteurs respectives $n+1$ et $n+2$, d'où notre
assertion. Supposons $n\geq2$ dans ce qui suit. Soit
$\mathfrak J=\mathfrak m\cap\mathfrak n=\mathfrak m\mathfrak n$, et
$R=1+\mathfrak J$ qui est une partie multiplicative de $B$~; si l'on pose
$B'=R^{-1}B$, l'idéal $\mathfrak J'=R^{-1}\mathfrak J$ est contenu dans le
radical de $B'$ (Bourbaki, \emph{Alg. comm.}, chap.~III, \S~3,
n\textsuperscript{o}~5, prop.~12)~; on sait que $X'=\operatorname{Spec}(B')$
s'identifie en tant qu'espace topologique à un sous-espace de
$X=\operatorname{Spec}(B)$, et qu'aux points $x$ de $X'$, les anneaux locaux
$\mathcal O_{X,x}$ et $\mathcal O_{X',x}$ sont les mêmes (\textbf{I}, 1.6.2).
Considérons alors dans $X'$ l'ensemble fermé $Y'=V(\mathfrak J')$, et soit
$S=X'-Y'$~; on sait (10.5.7) que $S$ est un préschéma de Jacobson, évidemment
irréductible et noethérien~; en outre les anneaux locaux
$\mathcal O_{S,s}=\mathcal O_{X,s}$ sont universellement caténaires pour tout $s\in S$
en vertu de (5.6.3), puisque $A$ est universellement caténaire (5.6.4). Pourtant,
il y a deux points fermés $a$, $b$ de $S$ tels que $\mathcal O_{S,a}$ et
$\mathcal O_{S,b}$ \emph{n'aient pas même dimension}, autrement dit $S$ ne vérifie
pas la condition $3^\circ$ de (10.6.1). Pour le voir, considérons les deux idéaux
maximaux $\mathfrak m'=R^{-1}\mathfrak m$, $\mathfrak n'=R^{-1}\mathfrak n$ de
$B'$, qui sont de hauteurs $n$ et $n+1$ respectivement~; on a
$\mathfrak J'=\mathfrak m'\cap\mathfrak n'$, et $\mathfrak J'$ n'est donc
contenu dans aucun idéal premier de $B'$ distinct de $\mathfrak m'$ et
$\mathfrak n'$, qui sont par suite les seuls idéaux maximaux de $B'$. Soient
$a'$, $b'$ les seuls points fermés de $X'$, correspondant à $\mathfrak m'$ et
$\mathfrak n'$. Il existe dans $B'$ un idéal premier non maximal
$\mathfrak p'\subset\mathfrak m'$ qui n'est pas contenu dans $\mathfrak n'$~:
il suffit de montrer qu'il y a dans $B$ un idéal premier non maximal contenu dans
$\mathfrak m$ et non dans $\mathfrak n$~; pour cela, on pourra par exemple
considérer la fibre de $\mathfrak m$ pour le morphisme correspondant à l'injection
$A[T_1]\to B=A[T_1,\ldots,T_n]$, et appliquer (6.1.2). En considérant une chaîne
maximale d'idéaux premiers entre $\mathfrak p'$ et $\mathfrak m'$ et remplaçant
$\mathfrak p'$ par l'avant-dernier idéal de cette chaîne, on peut donc supposer
que le point $a$ de $X'$ correspondant à $\mathfrak p'$ est tel que son adhérence
dans $X'$ soit $\{a,a'\}$~; comme $B'_{\mathfrak m'}$ est biéquidimensionnel,
$\mathfrak p'$ est alors de hauteur $n-1$. On construit de la même manière un
idéal premier non maximal $\mathfrak q'$ de $B'$ de hauteur $n$, tel que si $b$
est le point correspondant de $X'$, l'adhérence de $b$ dans $X'$ soit
$\{b,b'\}$. Cela étant, $a$ et $b$ sont dans $S$, donc \emph{fermés dans $S$},
et répondent par suite à la question.
\end{env}

\subsection{Profondeur rectifiée}
\label{subsection:IV.10.8-fr}

\begin{definition}[10.8.1]
\label{IV.10.8.1-fr}
Soient $X$ un préschéma localement noethérien, $\mathcal F$ un $\mathcal O_X$-Module
cohérent. Pour tout $x\in X$, on appelle \emph{profondeur rectifiée de $\mathcal F$
au point $x$} et on note $\operatorname{prof}^{*}_{x}(\mathcal F)$ le nombre (entier
$\geq0$ ou $+\infty$) égal à
\begin{equation}
\label{IV.10.8.1.1-fr}
\tag{10.8.1.1}
\operatorname{prof}^{*}_{x}(\mathcal F)
=\operatorname{prof}(\mathcal F_x)+\dim(\overline{\{x\}})
\end{equation}
où $\overline{\{x\}}$ est l'adhérence du point $x$ dans $X$. Pour toute partie
$Z$ de $X$, on appelle \emph{profondeur rectifiée de $\mathcal F$ le long de $Z$}
et on note $\operatorname{prof}^{*}_{Z}(\mathcal F)$ le nombre
\begin{equation}
\label{IV.10.8.1.2-fr}
\tag{10.8.1.2}
\operatorname{prof}^{*}_{Z}(\mathcal F)
=\inf_{x\in Z}\operatorname{prof}^{*}_{x}(\mathcal F).
\end{equation}
\end{definition}

En d'autres termes pour tout entier $n$, la relation
$\operatorname{prof}^{*}_{Z}(\mathcal F)\geq n$ équivaut à
$\operatorname{prof}^{*}_{x}(\mathcal F)\geq n$ pour tout $x\in Z$. Si $Z=X$, on
écrit $\operatorname{prof}^{*}(\mathcal F)$ au lieu de
$\operatorname{prof}^{*}_{X}(\mathcal F)$.

\oldpage[IV-3]{111}
\begin{remarks}[10.8.2]
\label{IV.10.8.2-fr}
(i) En tout point fermé $x\in X$, la profondeur rectifiée est égale à la
profondeur.

(ii) Soient $Y$ un sous-préschéma fermé de $X$, $j:Y\to X$ l'injection canonique,
$\mathcal G$ un $\mathcal O_Y$-Module cohérent. On sait (5.7.3, (vi)) que l'on a
$\operatorname{prof}(\mathcal G_x)=\operatorname{prof}(j_*(\mathcal G)_x)$ pour tout
$x\in Y$~; on en déduit que l'on a aussi
$\operatorname{prof}^{*}_{x}(\mathcal G)
=\operatorname{prof}^{*}_{x}(j_*(\mathcal G))$ pour tout $x\in Y$.

(iii) La notion de profondeur rectifiée n'a d'intérêt que lorsqu'elle est de
caractère \emph{local}, c'est-à-dire qu'elle ne change pas lorsqu'on remplace $X$
par un voisinage ouvert arbitraire $U$ de $x$. Cela exige évidemment que $x$ ne
soit pas isolé dans $\overline{\{x\}}$ lorsque $x$ n'est pas fermé et par suite
que $X$ soit un \emph{préschéma de Jacobson} (10.4.5.1)~; le plus souvent, il sera
aussi nécessaire de savoir que $\dim(U)=\dim(X)$ pour tout ouvert dense $U$ dans
$X$, et il faudra donc supposer que $X$ vérifie aussi les conditions $2^\circ$ et
$3^\circ$ de (10.6.1).
\end{remarks}

\begin{lemma}[10.8.3]
\label{IV.10.8.3-fr}
Soient $X$ un préschéma régulier et biéquidimensionnel, $\mathcal F$ un
$\mathcal O_X$-Module cohérent. Alors on a, pour tout $x\in X$,
\begin{equation}
\label{IV.10.8.3.1-fr}
\tag{10.8.3.1}
\operatorname{prof}^{*}_{x}(\mathcal F)
=\dim(X)-\dim.\operatorname{proj}(\mathcal F_x).
\end{equation}
\end{lemma}

En effet, comme $X$ est biéquidimensionnel, on a ($\mathbf 0$, 14.3.5.1)
\[
  \dim(\overline{\{x\}})=\dim(X)-\operatorname{codim}(\overline{\{x\}},X)
  =\dim(X)-\dim(\mathcal O_{X,x})
\]
en vertu de (5.1.2). D'autre part, comme $X$ est régulier, on a par
($\mathbf 0$, 17.3.4)
\[
  \operatorname{prof}(\mathcal F_x)
  =\dim(\mathcal O_{X,x})-\dim.\operatorname{proj}(\mathcal F_x)
\]
d'où le lemme.

\begin{corollary}[10.8.4]
\label{IV.10.8.4-fr}
Sous les hypothèses de (10.8.3), la fonction
$x\mapsto\operatorname{prof}^{*}_{x}(\mathcal F)$ est semi-continue inférieurement.
\end{corollary}

Cela résulte de (10.8.3.1), puisque
$x\mapsto\dim.\operatorname{proj}(\mathcal F_x)$ est semi-continue supérieurement
(6.11.1).

\begin{proposition}[10.8.5]
\label{IV.10.8.5-fr}
Soient $S$ un préschéma localement noethérien, $X$ un préschéma localement de
type fini sur $S$, $\mathcal F$ un $\mathcal O_X$-Module cohérent. On suppose que $S$
vérifie les conditions suivantes~: $1^\circ$ $S$ est un préschéma de Jacobson~;
$2^\circ$ $S$ est régulier~; $3^\circ$ les composantes irréductibles de $S$ sont
équicodimensionnelles. Alors la fonction
$x\mapsto\operatorname{prof}^{*}_{x}(\mathcal F)$ est semi-continue inférieurement
dans $X$~; autrement dit, pour tout entier $n$, l'ensemble $U_n$ des $x\in X$
tels que $\operatorname{prof}^{*}_{x}(\mathcal F)\geq n$ est ouvert.
\end{proposition}

Comme les anneaux locaux $\mathcal O_{S,s}$ de $S$ sont réguliers, ils sont
universellement caténaires (5.6.4)~; autrement dit, $S$ vérifie les conditions
$1^\circ$, $2^\circ$ et $3^\circ$ de (10.6.1), donc il en est de même de $X$
(10.6.1, (i)). La notion de profondeur rectifiée étant alors de caractère local
(10.8.2, (iii)), on peut se borner au cas où $S=\operatorname{Spec}(A)$ et
$X=\operatorname{Spec}(B)$ sont affines, $A$ étant un anneau régulier et $B$ une
$A$-algèbre de type fini, donc quotient d'un anneau de polynômes
$C=A[T_1,\ldots,T_n]$, et ce dernier est régulier ($\mathbf 0$, 17.3.7). On peut
donc supposer que $X$ est un sous-préschéma fermé d'un préschéma régulier $Y$
vérifiant également les conditions $1^\circ$, $2^\circ$ et $3^\circ$ de (10.6.1)~;
compte tenu de la remarque (10.8.2, (ii)), on est ainsi ramené au cas où $X$ est
en outre régulier et noethérien. Mais comme les anneaux locaux de $X$ sont alors
intègres, les composantes irréductibles de $X$ sont

\oldpage[IV-3]{112}
ouvertes (\textbf{I}, 6.1.10), et on peut par suite supposer aussi $X$
irréductible. Alors, comme les anneaux locaux de $X$ sont caténaires
($\mathbf 0$, 16.5.12), l'hypothèse $3^\circ$ de (10.6.1) entraîne que $X$ est
biéquidimensionnel ((5.1.5) et ($\mathbf 0$, 14.3.3))~; il suffit donc d'appliquer
(10.8.4).

On notera que si $S$ est le \emph{spectre d'un corps ou de $\mathbf Z$}, il
vérifie les conditions de (10.8.5).

\begin{corollary}[10.8.6]
\label{IV.10.8.6-fr}
Les hypothèses étant celles de (10.8.5), pour tout $x\in X$, le nombre
$\operatorname{prof}^{*}_{x}(\mathcal F)$ est l'unique entier $n$ ayant la propriété
suivante~: il existe un voisinage ouvert $U$ de $x$ dans $X$ tel que pour tout
point $x'\in U\cap\overline{\{x\}}$, fermé dans $U$, on ait
$\operatorname{prof}(\mathcal F_{x'})=n$. En particulier, pour que
$\operatorname{prof}^{*}_{x}(\mathcal F)\geq m$ (resp.
$\operatorname{prof}^{*}_{x}(\mathcal F)\leq m$), il faut et il suffit qu'il existe
un voisinage ouvert $V$ de $x$ dans $X$ tel que, pour tout
$x'\in V\cap\overline{\{x\}}$, fermé dans $V$, on ait
$\operatorname{prof}(\mathcal F_{x'})\geq m$ (resp.
$\operatorname{prof}(\mathcal F_{x'})\leq m$).
\end{corollary}

En effet, on peut se borner au cas où $x$ n'est pas fermé~; si
$\operatorname{prof}^{*}_{x}(\mathcal F)=n$, l'ensemble des $y\in X$ tels que
$\operatorname{prof}^{*}_{y}(\mathcal F)\leq n$ est fermé en vertu de (10.8.5), donc
contient $\overline{\{x\}}$, et en vertu de la semi-continuité inférieure de
$y\mapsto\operatorname{prof}^{*}_{y}(\mathcal F)$, il existe un voisinage ouvert $U$
de $x$ tel que $\operatorname{prof}^{*}_{y}(\mathcal F)=n$ pour tout
$y\in U\cap\overline{\{x\}}$, donc $\operatorname{prof}(\mathcal F_{x'})=n$ si
$x'\in U\cap\overline{\{x\}}$ est fermé dans $U$ (puisque la notion de profondeur
rectifiée est locale).

Pour les préschémas vérifiant les hypothèses de (10.8.5), la notion de profondeur
rectifiée peut donc se définir à l'aide des valeurs de la profondeur aux
\emph{points fermés} de $X$ (ces derniers formant un ensemble \emph{très dense}
dans toute partie fermée de $X$).

\begin{proposition}[10.8.7]
\label{IV.10.8.7-fr}
Soit $S$ un préschéma localement noethérien vérifiant les conditions $1^\circ$,
$2^\circ$ et $3^\circ$ de (10.6.1). Soient $X$ un préschéma localement de type
fini sur $S$, $\mathcal F$ un $\mathcal O_X$-Module cohérent,
$Y=\operatorname{Supp}(\mathcal F)$. Alors, pour tout $x\in Y$, on a
\begin{equation}
\label{IV.10.8.7.1-fr}
\tag{10.8.7.1}
\operatorname{prof}^{*}_{x}(\mathcal F)
=\dim_x(Y)-\operatorname{coprof}(\mathcal F_x).
\end{equation}
\end{proposition}

En effet, par définition, on a
$\operatorname{coprof}(\mathcal F_x)
=\dim(\mathcal F_x)-\operatorname{prof}(\mathcal F_x)$, et il résulte de (10.6.4) que
l'on a $\dim_x(Y)=\dim(\overline{\{x\}})+\dim(\mathcal F_x)$~; d'où (10.8.7.1) par
définition de $\operatorname{prof}^{*}_{x}(\mathcal F)$.

\begin{corollary}[10.8.8]
\label{IV.10.8.8-fr}
Les hypothèses sur $f$ et $\mathcal F$ étant celles de (9.9.1), la fonction
$x\mapsto\operatorname{prof}^{*}_{x}(\mathcal F_{f(x)})$ est constructible.
\end{corollary}

Notons que pour tout $s\in S$, la fibre $X_s$ étant un préschéma localement de
type fini sur $k(s)$, est un préschéma de Jacobson (10.4.7)~; comme en outre
$\operatorname{Spec}(k(s))$ vérifie les conditions de (10.6.1), on a
\[
  \operatorname{prof}^{*}_{x}(\mathcal F_{f(x)})
  =\dim_x(\operatorname{Supp}(\mathcal F_{f(x)}))
  -\operatorname{coprof}((\mathcal F_{f(x)})_x)
\]
(10.8.7). Or, si $Z=\operatorname{Supp}(\mathcal F)$, on a
$Z_{f(x)}=\operatorname{Supp}(\mathcal F_{f(x)})$ (\textbf{I}, 9.1.13) et $Z$ est
localement constructible (8.9.1)~; donc les fonctions
$x\mapsto\dim_x(\operatorname{Supp}(\mathcal F_{f(x)}))$ et
$x\mapsto\operatorname{coprof}((\mathcal F_{f(x)})_x)$ sont localement
constructibles ((9.9.1) et (9.9.3)), ce qui prouve la proposition.

\subsection{Spectres maximaux et ultrapréschémas}
\label{subsection:IV.10.9-fr}

Les résultats de ce numéro ne seront pas utilisés par la suite.

\begin{env}[10.9.1]
\label{IV.10.9.1-fr}
Soit $X$ un préschéma de Jacobson, et soit $S(X)$ l'\emph{espace annelé} dont
l'espace sous-jacent est le sous-espace des points fermés de $X$, et le faisceau
d'anneaux le faisceau induit sur ce sous-espace par $\mathcal O_X$, autrement dit
le faisceau d'anneaux $\psi^*(\mathcal O_X)$, en désignant par
$\psi:S(X)\to X$ l'injection canonique. Comme $\psi$ est un quasi-
\oldpage[IV-3]{113}
homéomorphisme, on a vu (10.2.8, (ii)) que si
$\theta:\mathcal O_X\to\psi_*(\mathcal O_{S(X)})$ est l'homomorphisme de
faisceaux d'anneaux tel que
$\theta^\sharp:\psi^*(\mathcal O_X)\to\mathcal O_{S(X)}$ soit l'identité,
alors $j_X=(\psi,\theta)$ est un \emph{quasi-isomorphisme} d'espaces annelés,
et $\mathcal F\mapsto j_X^*(\mathcal F)$ une équivalence de la catégorie des
$\mathcal O_X$-Modules et de celle des $\mathcal O_{S(X)}$-Modules. Il est clair
que dans cette équivalence, aux $\mathcal O_X$-Modules localement libres (resp.
cohérents) correspondent les $\mathcal O_{S(X)}$-Modules localement libres
(resp. cohérents)~; en outre, si $\mathcal L$ est un $\mathcal O_X$-Module
localement libre et si $U$ est un ouvert de $X$ tel que $\mathcal L|U$ soit
isomorphe à $(\mathcal O_X|U)^n$, $j_X^*(\mathcal L)$ est tel que
$j_X^*(\mathcal L)|(U\cap S(X))$ soit isomorphe à
$(\mathcal O_{S(X)}|(U\cap S(X)))^n$.
\end{env}

\begin{env}[10.9.2]
\label{IV.10.9.2-fr}
Soient $X$, $Y$ deux préschémas de Jacobson et
$f=(\rho,\lambda):X\to Y$ un morphisme \emph{localement de type fini}. On a vu
(10.4.7) que l'on a $\rho(S(X))\subset S(Y)$ et par restriction de $\rho$ à
$S(X)$, on définit donc une application continue $S(\rho):S(X)\to S(Y)$.
D'autre part, pour tout ouvert $V$ de $Y$, on définit par composition un
homomorphisme d'anneaux
\[
  \Gamma(V\cap S(Y),\mathcal O_{S(Y)})\longrightarrow
  \Gamma(V,\mathcal O_Y)\xrightarrow{\lambda_V}
  \Gamma(f^{-1}(V),\mathcal O_X)\longrightarrow
  \Gamma(f^{-1}(V)\cap S(X),\mathcal O_{S(X)})
\]
où les deux isomorphismes extrêmes ont été définis dans (10.9.1)~; il est clair
que cela définit un homomorphisme de faisceaux d'anneaux
$S(\lambda):\mathcal O_{S(Y)}\to\rho_*(\mathcal O_{S(X)})$ (en se rappelant
que les ouverts de $S(X)$ (resp. $S(Y)$) correspondent biunivoquement à ceux de
$X$ (resp. $Y$) (10.2.1))~; on obtient ainsi un morphisme d'espaces annelés
\[
  S(f)=(S(\rho),S(\lambda)):(S(X),\mathcal O_{S(X)})\longrightarrow
  (S(Y),\mathcal O_{S(Y)})
\]
tel que le diagramme
\[
\begin{tikzcd}
S(X) \ar[r,"S(f)"] \ar[d,"j_X"'] & S(Y) \ar[d,"j_Y"] \\
X \ar[r,"f"'] & Y
\end{tikzcd}
\]
soit commutatif~; en outre, si $Z$ est un troisième préschéma de Jacobson et
$g:Y\to Z$ un morphisme localement de type fini, il est clair que
$S(g\circ f)=S(g)\circ S(f)$. On a ainsi défini un \emph{foncteur covariant}
$S:\mathbf C\to\mathbf C'$, où $\mathbf C'$ est la catégorie des \emph{espaces
annelés en anneaux locaux}, et $\mathbf C$ la catégorie dont les objets sont les
\emph{préschémas de Jacobson} et les morphismes sont les morphismes
\emph{localement de type fini} entre préschémas de Jacobson.
\end{env}

\begin{env}[10.9.3]
\label{IV.10.9.3-fr}
Proposons-nous de déterminer la sous-catégorie $\mathbf C''$ de $\mathbf C'$
formée des espaces annelés isomorphes aux $S(X)$ et dont les morphismes
proviennent des $S(f)$. Supposons d'abord que $X=\operatorname{Spec}(A)$, où
$A$ est un anneau de Jacobson~; alors $S(X)$ est l'ensemble des \emph{idéaux
maximaux} de $A$, muni~: $1^\circ$ de la topologie induite par celle de $X$, de
sorte qu'une base de cette topologie est formée des
$D^m(h)=D(h)\cap S(X)$, ensemble des idéaux maximaux $\mathfrak m$ de $A$ tels
que $h\notin\mathfrak m$, où $h$ parcourt $A$~; $2^\circ$ du faisceau d'anneaux
$\mathcal O_{S(X)}$ tel que
$\Gamma(D^m(h),\mathcal O_{S(X)})=A_h$. Nous dirons que cet espace annelé est
le \emph{spectre maximal} de l'anneau de Jacobson $A$ et nous le noterons
$\operatorname{Spm}(A)$.

On notera que si $j:D(h)\to X$ est l'injection canonique, l'espace annelé
induit sur $D^m(h)$ par $S(X)$ est $S(D(h))$ et l'injection canonique
$D^m(h)\to S(X)$ d'espaces annelés est égale à $S(j)$.

Soient $B$ un second anneau de Jacobson, $Y=\operatorname{Spec}(B)$,
$\varphi:B\to A$ un homomorphisme d'anneaux faisant de $A$ une $B$-algèbre de
type fini, $f=({}^a\varphi,\widetilde\varphi):X\to Y$ le morphisme
correspondant de préschémas, et
\[
  S(f):\operatorname{Spm}(A)\longrightarrow\operatorname{Spm}(B)
\]
le morphisme d'espaces annelés correspondant à $f$. Il est clair que
$S(f)=(\psi,\theta)$ est un morphisme d'espaces annelés en anneaux locaux,
c'est-à-dire ($\mathrm{Err}_{\mathrm{II}}$, 1.8.2) que pour tout
$x\in\operatorname{Spm}(A)$, $\theta_x^\sharp$ est un homomorphisme local.
Réciproquement~:
\end{env}

\begin{proposition}[10.9.4]
\label{IV.10.9.4-fr}
Soient $A$, $B$ deux anneaux de Jacobson. Si
$u=(\psi,\theta):\operatorname{Spm}(A)\to\operatorname{Spm}(B)$ est un
morphisme d'espaces annelés en anneaux locaux tel que
$\Gamma(\theta):B\to A$ fasse de $A$ une $B$-algèbre de type fini, il existe
un morphisme de préschémas
$f:\operatorname{Spec}(A)\to\operatorname{Spec}(B)$ et un seul tel que
$u=S(f)$.
\end{proposition}

L'unicité de $f$ est évidente, puisque si
$f=({}^a\varphi,\widetilde\varphi)$, on doit avoir
$\varphi=\Gamma(\theta)$~; il reste à voir que $S(f)$ est défini et qu'on a
bien $u=S(f)$. Or, la première assertion résulte de ce que $\varphi$ est
supposé faire de $A$ une $B$-algèbre de type fini, et par suite
$f(\operatorname{Spm}(A))\subset\operatorname{Spm}(B)$~; le fait que
$\theta_x^\sharp$ est un homomorphisme local pour tout
$x\in\operatorname{Spm}(A)$ permet alors de répéter le raisonnement de
(\textbf I, 1.7.3) en se bornant aux points $x$ de $\operatorname{Spm}(A)$~:
on montre ainsi successivement que $\psi(x)={}^a\varphi(x)$ pour tout
$x\in\operatorname{Spm}(A)$, puis que
$\theta_x^\sharp=\varphi_x:B_{\psi(x)}\to A_x$, ce qui achève de prouver que
$u=S(f)$.

\oldpage[IV-3]{114}

\begin{env}[10.9.5]
\label{IV.10.9.5-fr}
Considérons maintenant un espace annelé $(X,\mathcal O_X)$~; nous dirons qu'une
partie ouverte $U$ de $X$ est \emph{ultra-affine} si l'espace annelé induit
$(U,\mathcal O_X|U)$ est isomorphe à un spectre maximal
$\operatorname{Spm}(A)$, où $A$ est un anneau de Jacobson. Nous dirons que $X$
est un \emph{ultra-préschéma} si tout point de $X$ admet un voisinage ouvert
ultra-affine. On montre, comme dans (\textbf I, 2.1.3 et 2.1.4) que les ouverts
ultra-affines forment une \emph{base} de la topologie de $X$ et que $X$ est un
espace de Kolmogoroff. Si $Y$ est un second ultra-préschéma, nous dirons qu'un
morphisme d'espaces annelés $f:X\to Y$ est un \emph{morphisme
d'ultra-préschémas} s'il vérifie les conditions suivantes~: $1^\circ$ $f$ est
un morphisme d'espaces annelés en anneaux locaux~; $2^\circ$ pour tout $x\in X$,
il y a un voisinage ouvert ultra-affine $V$ de $f(x)$ dans $Y$ et un voisinage
ouvert ultra-affine $U$ de $x$ dans $X$ tels que $f(U)\subset V$ et que
l'homomorphisme
$\Gamma(V,\mathcal O_Y)\to\Gamma(U,\mathcal O_X)$ correspondant à $f$ fasse
de $\Gamma(U,\mathcal O_X)$ une $\Gamma(V,\mathcal O_Y)$-algèbre de type fini.

Il est immédiat qu'on définit bien ainsi des morphismes, le composé de deux
morphismes en étant un troisième grâce à la remarque finale de (10.9.3). Il
est clair que la catégorie $\mathbf C''_0$ ainsi définie est une sous-catégorie
de $\mathbf C'$ qui contient $\mathbf C''$~; nous nous proposons de montrer que
$\mathbf C''=\mathbf C''_0$, autrement dit~:
\end{env}

\begin{proposition}[10.9.6]
\label{IV.10.9.6-fr}
Le foncteur $X\mapsto S(X)$ de $\mathbf C$ dans $\mathbf C''_0$ est une
équivalence de catégories.
\end{proposition}

$1^\circ$ Montrons d'abord que le foncteur $X\mapsto S(X)$ est
\emph{pleinement fidèle}, autrement dit que pour $X$, $Y$ dans $\mathbf C$,
l'application canonique
\[
  \operatorname{Hom}_{\mathbf C}(X,Y)\longrightarrow
  \operatorname{Hom}_{\mathbf C''_0}(S(X),S(Y))
\]
est \emph{bijective}. En premier lieu elle est \emph{injective}~: soient en
effet $f$, $g$ deux morphismes localement de type fini de $X$ dans $Y$ et
supposons que $S(f)=S(g)$. Cela entraîne d'abord que pour tout ouvert $V$ de
$Y$, on a
$f^{-1}(V)\cap S(X)=g^{-1}(V)\cap S(X)$, donc (10.3.1 et 10.2.7)
$f^{-1}(V)=g^{-1}(V)$~; il suffit donc de prouver que pour tout ouvert affine
$V$ de $Y$, $f$ et $g$ coïncident dans
$f^{-1}(V)=g^{-1}(V)$, autrement dit, on est ramené au cas où
$Y=\operatorname{Spec}(B)$ est le spectre d'un anneau de Jacobson $B$. Il
suffit (\textbf I, 2.2.4) de montrer que les homomorphismes d'anneaux
$B\to\Gamma(X,\mathcal O_X)$ correspondant à $f$ et $g$ sont alors les mêmes.
Or, pour tout $s\in B$, les images de $s$ par ces homomorphismes sont deux
sections de $\mathcal O_X$ au-dessus de $X$ qui, par hypothèse, induisent la
même section au-dessus de $S(X)$~; on sait donc (10.2.8) que ces sections sont
identiques, d'où notre assertion. Prouvons en second lieu que tout morphisme
$h:S(X)\to S(Y)$ (pour la catégorie $\mathbf C''_0$) est de la forme $S(f)$,
où $f:X\to Y$ est un morphisme localement de type fini. Par hypothèse, il
existe un recouvrement ouvert ultra-affine $(U'_\lambda)$ (resp.
$(V'_\mu)$) de $S(X)$ (resp. $S(Y)$) tel que pour tout $\lambda$,
$h(U'_\lambda)$ soit contenu dans un $V'_\mu$ et qu'il corresponde au
morphisme $h_\lambda:U'_\lambda\to V'_\mu$, restriction de $h$, un
homomorphisme d'anneaux
$\Gamma(V'_\mu,\mathcal O_{S(Y)})\to
\Gamma(U'_\lambda,\mathcal O_{S(X)})$ faisant du second anneau une algèbre de
type fini sur le premier. On peut supposer que
$U'_\lambda=U_\lambda\cap S(X)$ et $V'_\mu=V_\mu\cap S(Y)$, où
$U_\lambda$ et $V_\mu$ sont des ouverts affines uniquement déterminés dans
$X$ et $Y$ respectivement~; si l'on prouve que, pour tout $\lambda$,
$h_\lambda=S(f_\lambda)$ où $f_\lambda:U_\lambda\to V_\mu$ est un morphisme
de type fini, alors il résulte de la première partie du raisonnement,
appliquée aux restrictions de $f_\alpha$ et $f_\beta$ à
$U_\alpha\cap U_\beta$, que les $f_\lambda$ sont les restrictions d'un même
morphisme $f:X\to Y$, et on aura évidemment $h=S(f)$. On est ainsi ramené au
cas où $X$ et $Y$ sont affines, et la conclusion résulte alors de (10.9.4).

$2^\circ$ Il reste à prouver que tout ultra-préschéma $X'$ est de la forme
$S(X)$ pour un préschéma de Jacobson $X$ (qui sera nécessairement unique à
isomorphisme près, en vertu de $1^\circ$). Il y a un recouvrement
$(U'_\alpha)$ de $X'$ par des ouverts ultra-affines, dont chacun est de la
forme $S(U_\alpha)$, $U_\alpha$ étant le spectre d'un anneau de Jacobson. Pour
tout couple d'indices $\alpha$, $\beta$, considérons l'unique ouvert
$U_{\alpha\beta}$ de $U_\alpha$ dont la trace sur $U'_\alpha$ est
$U'_\alpha\cap U'_\beta$~; en vertu de $1^\circ$, l'automorphisme identique de
$U'_\alpha\cap U'_\beta$ est de la forme $S(\theta_{\alpha\beta})$, où
$\theta_{\alpha\beta}:U_{\alpha\beta}\to U_{\beta\alpha}$ est un
isomorphisme de préschémas. On constate aussitôt (en vertu du $1^\circ$) que la
famille $(\theta_{\alpha\beta})$ vérifie la condition de recollement
($\mathbf 0_{\mathrm I}$, 4.1.7), et que cette famille définit donc un
préschéma $X$, tel que les $U_\alpha$ s'identifient à des ouverts affines de
$X$~; il est clair alors que l'on a $X'=S(X)$, ce qui achève la démonstration.

\subsection{Espaces algébriques de Serre}
\label{subsection:IV.10.10-fr}

\begin{env}[10.10.1]
\label{IV.10.10.1-fr}
Le langage introduit par Serre dans (FAC) est parfois commode, notamment dans
les questions où l'intérêt majeur s'attache aux points rationnels sur le corps
de base $k$ (algébriquement clos par hypothèse) des «~variétés algébriques~»
sur $k$ que l'on considère. Nous allons esquisser ici ce langage en le
rattachant aux considérations qui précèdent, pour permettre au lecteur de
traduire les énoncés de Serre en langage des schémas. Il est d'ailleurs
possible de développer le langage de Serre également pour des préschémas sur
un corps non algébriquement clos (et même sur un anneau artinien)~; mais cela
introduit des complications techniques considérables, et d'ailleurs, sur un
corps de base arbitraire, les avantages (surtout psychologiques) du point de
vue de Serre disparaissent~; aussi
\oldpage[IV-3]{115}
nous tiendrons-nous dans le cadre fixé par Serre. Le présent numéro, tout comme
le précédent, ne sera d'ailleurs pas utilisé dans la suite de ce Traité, et
nous nous bornerons donc à de brèves indications.
\end{env}

\begin{env}[10.10.2]
\label{IV.10.10.2-fr}
Étant donné un ultra-préschéma fixe $R$, on peut naturellement (comme dans
toute catégorie) définir la notion de $R$-\emph{ultra-préschéma}. Considérons
en particulier un corps algébriquement clos $k$~; $\operatorname{Spm}(k)$ est
alors identique à $\operatorname{Spec}(k)$~; nous dirons qu'un
$\operatorname{Spm}(k)$-ultra-préschéma $X'$ est un \emph{espace
préalgébrique sur $k$}~: c'est donc un espace $k$-annelé en anneaux locaux dont
tout point a un voisinage ouvert isomorphe au spectre maximal d'une $k$-algèbre
de type fini~; il revient au même de dire (par (10.9.6)) que $X'=S(X)$, où $X$
est un préschéma localement de type fini sur $k$. Si $X$, $Y$, $Z$ sont trois
$k$-préschémas localement de type fini, il en est de même de
$X\times_ZY$ (\textbf I, 3.4), donc en vertu de (10.9.6), la notion de produit
existe dans la catégorie des espaces préalgébriques sur $k$ (aussi bien le
produit «~sur $k$~» $X'\times_kY'$ que le «~produit fibré~»
$X'\times_{Z'}Y'$, où $X'$, $Y'$, $Z'$ sont trois espaces préalgébriques sur
$k$). On peut donc définir le morphisme diagonal
$\Delta_{X'}:X'\to X'\times_kX'$ (qui n'est autre d'ailleurs que
$S(\Delta_X)$)~; on dit que $X'$ est un \emph{espace algébrique sur $k$} si
$X$ est un schéma, et il revient au même de dire que l'image de $\Delta_{X'}$
est une partie \emph{fermée} de $X'\times_kX'$.
\end{env}

\begin{env}[10.10.3]
\label{IV.10.10.3-fr}
Les simplifications qui proviennent de l'hypothèse que $k$ est algébriquement
clos tiennent d'abord à ce que, pour un $k$-préschéma $X$, localement de type
fini, il y a correspondance biunivoque entre \emph{points fermés de $X$},
\emph{points de $X$ à valeurs dans $k$} (\textbf I, 3.4.4) et \emph{points de
$X$ rationnels sur $k$} (\textbf I, 3.4.5), en vertu de (\textbf I, 6.4.2).
Cela montre en particulier (en vertu de (\textbf I, 3.4.3.1)) que pour deux
espaces $k$-préalgébriques $X'$, $Y'$, l'ensemble sous-jacent au produit
$X'\times_kY'$ est identique à l'\emph{ensemble produit} $X'\times Y'$ des
ensembles sous-jacents (mais bien entendu la topologie de l'espace sous-jacent
à $X'\times_kY'$ n'est pas la \emph{topologie produit} des topologies des
espaces sous-jacents à $X'$ et $Y'$, elle est en général strictement plus fine
que cette dernière).

D'autre part, les anneaux locaux $\mathcal O_x$ aux points d'un espace
préalgébrique $X'$ sur $k$ sont des $k$-algèbres dont on vient de voir que le
corps résiduel est \emph{isomorphe à $k$}~; si $A$ et $B$ sont deux telles
$k$-algèbres locales, tout $k$-homomorphisme $\varphi:A\to B$ est
nécessairement \emph{local}~: en effet, si un élément $x$ de l'idéal maximal de
$A$ était tel que $\varphi(x)$ soit inversible, il existerait $\lambda\in k$
non nul tel que $\varphi(x-\lambda.1)$ appartienne à l'idéal maximal de $B$, ce
qui est absurde puisque $x-\lambda.1$ est inversible dans $A$. On conclut
aussitôt de là que si $X'=S(X)$, $Y'=S(Y)$ sont deux espaces préalgébriques sur
$k$, tout morphisme $X'\to Y'$ d'espaces $k$-annelés est aussi un morphisme
d'espaces $k$-annelés en anneaux locaux (\textbf I, 1.8.2~; cf.
$\mathrm{Err}_{\mathrm{II}}$). D'ailleurs, avec les notations précédentes, si
$A$ et $B$ sont des $k$-algèbres de type fini, $\varphi$ fait de $B$ une
$A$-algèbre de type fini~; donc tout morphisme $X'\to Y'$ de $k$-espaces
annelés est, en vertu de (10.9.6), de la forme $S(f)$, où $f:X\to Y$ est un
morphisme de $k$-préschémas.

Enfin, pour tout ouvert $U$ de $X'$, toute section
$s\in\Gamma(U,\mathcal O_{X'})$, et tout $x\in U$, $s(x)$
($\mathbf 0_{\mathrm I}$, 5.5.1) s'identifie à un élément de $k$, et l'on a
ainsi associé à $s$ une application $\overline s:x\mapsto s(x)$ de $U$ dans
$k$, autrement dit une section au-dessus de $U$ du faisceau $\mathcal A(X')$
des germes d'applications de $X'$ dans $k$~; comme l'application
$h_U:s\mapsto\overline s$ est évidemment un homomorphisme d'anneaux
$\Gamma(U,\mathcal O_{X'})\to\Gamma(U,\mathcal A(X'))$ et commute aux
restrictions à un ouvert $V\subset U$, les $h_U$ définissent un homomorphisme
de faisceaux d'anneaux $h:\mathcal O_{X'}\to\mathcal A(X')$. Si l'on prend
pour $U$ un ouvert ultra-affine $\operatorname{Spm}(A)$, où $A$ est un anneau
de Jacobson, dire que $\overline s=0$ signifie que pour tout idéal maximal
$\mathfrak m$ de $A$, $s$ appartient à $\mathfrak m$, ou encore que $s$ est
dans le \emph{radical} de $A$~; mais comme $A$ est un anneau de Jacobson, son
radical est égal à son nilradical~; pour que $h_U$ soit injective, il faut et
il suffit par suite que $A$ soit réduit.
\end{env}

\begin{env}[10.10.4]
\label{IV.10.10.4-fr}
On dit que l'espace $k$-préalgébrique $X'=S(X)$ est \emph{réduit} s'il en est
ainsi de $X$~; comme l'ensemble des points $x\in X$ où $X$ est réduit est
ouvert ($\mathbf 0_{\mathrm I}$, 5.2.2), son complémentaire contient au moins
un point fermé s'il est non vide (5.1.11), et il revient donc au même de dire
que $X'$ est réduit ou que chacun de ses anneaux locaux $\mathcal O_x$ (pour
$x\in X'$) est réduit. On vient de voir dans (10.10.3) que pour que
l'homomorphisme $h:\mathcal O_{X'}\to\mathcal A(X')$ soit \emph{injectif}, il
faut et il suffit que $X'$ soit réduit. Dans (FAC), Serre se borne en fait aux
espaces préalgébriques \emph{réduits}, ce qui lui permet de définir
$\mathcal O_{X'}$ comme un \emph{sous-faisceau} de $\mathcal A(X')$. On notera
que si $X'$ et $Y'$ sont des $k$-espaces préalgébriques réduits, il en est de
même de $X'\times_kY'$~: en effet, tout revient à voir que si $A$ et $B$ sont
deux $k$-algèbres de type fini réduites, il en est de même de $A\otimes_kB$~;
mais nous avons vu que les radicaux de $A$ et $B$ sont alors réduits à $0$, et
comme $k$ est algébriquement clos, $A$ et $B$ sont des algèbres «~séparables~»
sur $k$ au sens de Bourbaki (Bourbaki, \emph{Alg.}, chap.~VIII, \S~7,
n$^\circ$~5, prop.~5)~; donc $A\otimes_kB$ est sans radical (\emph{loc.
cit.}, n$^\circ$~6, cor.~3 du th.~3), et puisque c'est un anneau de Jacobson,
il est réduit. Toutefois, si $Z'$ est un troisième espace préalgébrique sur
$k$, le «~produit fibré~» $X'\times_{Z'}Y'$ de deux espaces préalgébriques
réduits sur $Z'$ n'est pas réduit en général, ce qui implique que la catégorie
de ces espaces est insuffisante dans de nombreuses questions (notamment dans
la théorie des groupes algébriques). Mais comme on l'a vu ci-dessus, on peut
garder le langage de Serre sans se limiter, comme ce dernier (qui en outre ne
considère que des espaces préalgébriques quasi-compacts), au cas des espaces
préalgébriques réduits.
\end{env}

\oldpage[IV-3]{116}

\begin{env}[10.10.5]
\label{IV.10.10.5-fr}
Enfin, on peut aussi considérer des ultra-préschémas sur un corps $k$
quelconque en conservant un langage qui reste proche de celui de Serre, et en
introduisant, comme chez Weil, une extension algébriquement close fixe $K$ de
$k$ (choisie assez grande, par exemple de degré de transcendance infini sur
$k$, pour avoir suffisamment de «~points génériques~» au sens de Weil). À tout
préschéma $X$ localement de type fini sur $k$, on associe alors l'ensemble
$S_K(X)=S(X\otimes_kK)$ des points de $X$ à valeurs dans $K$~; on a une
application canonique $j:S_K(X)\to X$ que l'on démontre être un
quasi-homéomorphisme quand on munit $S_K(X)$ de la topologie image réciproque
de celle de $X$ par $j$~; on munit $S_K(X)$ du faisceau de $k$-algèbres
$j^*(\mathcal O_X)$, et on obtient ainsi une sous-catégorie de la catégorie des
espaces $k$-annelés en anneaux locaux, que l'on pourrait appeler catégorie des
$(k,K)$-espaces préalgébriques. On peut montrer qu'on y peut encore définir des
produits et généraliser les résultats de (10.10.3) et (10.10.4)
($\mathcal A(X')$ étant ici remplacé par le faisceau $\mathcal A_K(X')$ des
germes d'applications de $X'$ dans $K$). Toutefois ce point de vue fait jouer
un rôle artificiel à un surcorps arbitrairement choisi de $k$, et nous ne le
signalons que pour le rejeter.
\end{env}
